\documentclass[12pt,oneside]{book}

%============================================================
% PACKAGES
%============================================================

%--------------------
% Page + font
%--------------------
\usepackage[margin=1in]{geometry}
\usepackage[T1]{fontenc}
\usepackage[utf8]{inputenc}
\usepackage{lmodern}
\usepackage{microtype}
\usepackage{parskip}


%--------------------
% Math
%--------------------
\usepackage{amsmath, amssymb, amsfonts, amsthm, mathtools}
\usepackage{bm}
\usepackage{bbm}
\usepackage{mathrsfs}
\usepackage{dsfont}

%--------------------
% Tables + figures
%--------------------
\usepackage{graphicx}
\usepackage{float}
\usepackage{booktabs}
\usepackage{array}
\usepackage{xltabular}
\usepackage{multirow}
\usepackage{caption}
\usepackage{subcaption}

%--------------------
% Lists
%--------------------
\usepackage{enumitem}
\setlist[itemize]{topsep=2pt,itemsep=2pt}
\setlist[enumerate]{topsep=2pt,itemsep=2pt}

%--------------------
% Colors + hyperlinks
%--------------------
\usepackage[dvipsnames]{xcolor}
\usepackage[
    colorlinks=true,
    linkcolor=MidnightBlue,
    citecolor=MidnightBlue,
    urlcolor=MidnightBlue
]{hyperref}
\usepackage[nameinlink,capitalise]{cleveref}

%--------------------
% Algorithms
%--------------------
\usepackage{algorithm}
\usepackage{algpseudocode}

%--------------------
% Code blocks
%--------------------
\usepackage{listings}

\lstdefinestyle{codestyle}{
    basicstyle=\ttfamily\small,
    breaklines=true,
    frame=single,
    numbers=left,
    numberstyle=\tiny,
    keywordstyle=\color{MidnightBlue},
    commentstyle=\color{ForestGreen},
    stringstyle=\color{BrickRed},
    showstringspaces=false,
    tabsize=4
}

\lstset{style=codestyle}

%--------------------
% Section styling
%--------------------
\usepackage{titlesec}
\usepackage{tocloft}

\setcounter{tocdepth}{2}
\setcounter{secnumdepth}{3}

%============================================================
% THEOREMS
%============================================================

\numberwithin{equation}{section}

\theoremstyle{definition}
\newtheorem{definition}{Definition}[section]
\newtheorem{example}{Example}[section]
\newtheorem{exercise}{Exercise}[section]
\newtheorem{remark}{Remark}[section]

\theoremstyle{plain}
\newtheorem{theorem}{Theorem}[section]
\newtheorem{lemma}{Lemma}[section]
\newtheorem{proposition}{Proposition}[section]
\newtheorem{corollary}{Corollary}[section]

\theoremstyle{remark}
\newtheorem*{note}{Note}

%============================================================
% CUSTOM COMMANDS
%============================================================

% Sets
\newcommand{\R}{\mathbb{R}}
\newcommand{\N}{\mathbb{N}}
\newcommand{\Z}{\mathbb{Z}}
\newcommand{\Q}{\mathbb{Q}}
\newcommand{\C}{\mathbb{C}}

% Probability
\newcommand{\E}{\mathbb{E}}
\newcommand{\Var}{\mathrm{Var}}
\newcommand{\Cov}{\mathrm{Cov}}
\newcommand{\Corr}{\mathrm{Corr}}
\newcommand{\Prob}{\mathbb{P}}

\newcommand{\ind}{\mathbbm{1}}
\newcommand{\iid}{\overset{\text{iid}}{\sim}}
\newcommand{\convd}{\overset{d}{\to}}
\newcommand{\convp}{\overset{p}{\to}}
\newcommand{\as}{\overset{a.s.}{\to}}

% Operators
\newcommand{\given}{\,\middle|\,}
\newcommand{\argmin}{\operatorname*{arg\,min}}
\newcommand{\argmax}{\operatorname*{arg\,max}}

% Linear algebra / analysis
\newcommand{\norm}[1]{\left\lVert #1 \right\rVert}
\newcommand{\abs}[1]{\left\lvert #1 \right\rvert}
\newcommand{\inner}[2]{\left\langle #1, #2 \right\rangle}

\newcommand{\dd}{\,\mathrm{d}}
\newcommand{\grad}{\nabla}
\newcommand{\Hess}{\nabla^2}

\newcommand{\tr}{\operatorname{tr}}
\newcommand{\rank}{\operatorname{rank}}
\newcommand{\diag}{\operatorname{diag}}

% Distributions
\newcommand{\Normal}{\mathcal{N}}
\newcommand{\Bern}{\mathrm{Bernoulli}}
\newcommand{\Bin}{\mathrm{Binomial}}
\newcommand{\Pois}{\mathrm{Poisson}}
\newcommand{\Gam}{\mathrm{Gamma}}
\newcommand{\BetaDist}{\mathrm{Beta}}
\newcommand{\Exp}{\mathrm{Exponential}}
\newcommand{\Dir}{\mathrm{Dirichlet}}
\newcommand{\InvGam}{\mathrm{Inv\text{-}Gamma}}

% Title block
\title{EN.553.730 Statistical Theory}
\author{Jonathan Y. Ma}
\date{June 2026}

\begin{document}

\maketitle

\noindent \emph{This course was taught by Dr. Alden Green in the Fall 2026 Semester. The notes are transcribed by Jonathan Ma and may contain errors and omitted content, so any issues are to be directed towards me.}

\tableofcontents

\part{Data Reduction}

\section{Introduction}

Statistical theory begins with an experiment whose outcome is random but whose
distribution is only partially known.  Let
\((\mathcal X,\mathcal A)\) be the sample space and let
\[
    \mathcal P=\{P_\theta:\theta\in\Theta\}
\]
be a family of probability measures on it.  The pair
\((\mathcal X,\mathcal P)\), with the underlying $\sigma$-field understood, is
called a \emph{statistical experiment}.  We observe $X\sim P_\theta$ and wish
to learn something about the unknown parameter $\theta$.

The raw observation often contains more detail than is relevant to this goal.
A \emph{statistic} is a measurable map
\[
    T:\mathcal X\longrightarrow\mathcal T.
\]
Replacing $X$ by $T(X)$ is a data reduction.  Such a reduction is attractive
when $T$ is simpler than $X$, but simplicity alone is not enough: we would
like to discard no information about $\theta$.

\begin{example}[Repeated coin flips]
Let $X_1,\ldots,X_n\iid\Bern(p)$, where $p\in[0,1]$.  The full data are the
binary string $X=(X_1,\ldots,X_n)\in\{0,1\}^n$.  If the order of the flips is
irrelevant, a natural reduction is
\[
    S=\sum_{i=1}^n X_i,
\]
the total number of heads.  The sample space has been reduced from $2^n$
possible strings to $n+1$ possible counts.  We will see that this reduction
loses no information about $p$.
\end{example}

There are two complementary ways to make ``no loss of information'' precise.
The decision-theoretic formulation asks whether every procedure based on $X$
can be reproduced, with the same performance, using only $T(X)$.  The
probabilistic formulation asks whether, after $T(X)$ is known, the remaining
randomness in $X$ has a distribution that is independent of $\theta$.  The
notion of sufficiency connects these two points of view.

\section{The Decision-Theoretic Framework}

A statistical problem requires more than a model.  Let $\mathcal D$ be an
action space, and let
\[
    L:\Theta\times\mathcal D\longrightarrow[0,\infty]
\]
be a loss function.  The value $L(\theta,a)$ is the cost of taking action $a$
when the true parameter is $\theta$.  A nonrandomized \emph{decision rule} is
a measurable function $\delta:\mathcal X\to\mathcal D$.  Its frequentist risk
is
\begin{equation}
    R(\theta,\delta)
    =\E_\theta\!\left[L\bigl(\theta,\delta(X)\bigr)\right].
    \label{eq:risk-definition}
\end{equation}
Thus a rule is evaluated not by a single number, but by the function
$\theta\mapsto R(\theta,\delta)$.

Common examples include the following.
\begin{itemize}
    \item For point estimation of $g(\theta)\in\R$, squared-error loss is
    $L(\theta,a)=(a-g(\theta))^2$.
    \item Absolute-error loss is $L(\theta,a)=\abs{a-g(\theta)}$.
    \item For testing $H_0:\theta\in\Theta_0$ against
    $H_1:\theta\in\Theta_1$, the action space is $\{0,1\}$, where action $1$
    means rejection.  Zero-one loss assigns a cost to an incorrect decision;
    asymmetric loss may assign different costs to the two kinds of error.
\end{itemize}

Randomization is occasionally useful.  Formally, a randomized rule is a
Markov kernel $Q(\dd a\mid x)$ from $\mathcal X$ to $\mathcal D$: after
observing $x$, the action is drawn from $Q(\cdot\mid x)$.  A deterministic
rule is the special case $Q(\dd a\mid x)=\delta_{\delta(x)}(\dd a)$.

Suppose only $T=T(X)$ is retained.  Then a rule based on the reduced data has
the form $d(T)$, or more generally is a kernel from $\mathcal T$ to
$\mathcal D$.  The central question is whether this restriction makes any
decision problem harder.  If the conditional distribution of $X$ given $T$
does not depend on $\theta$, a rule based on $X$ can be simulated from $T$:
given $T=t$, generate $X^*$ from that conditional distribution and apply the
original rule to $X^*$.  Under every $P_\theta$, $X^*$ has the same marginal
law as $X$, so the simulated rule has exactly the same risk function.

This observation motivates sufficiency.  It also makes clear why randomized
rules are natural in the most general decision-theoretic statement: the
reconstruction of $X$ from $T$ may itself require randomization.  For convex
losses and point estimation, conditional expectation will often remove that
randomization and improve risk; this is the content of the Rao--Blackwell
theorem.

\begin{example}[Estimating a coin bias]
For $X_1,\ldots,X_n\iid\Bern(p)$ and squared-error loss, consider the rule
$\delta(X)=X_1$.  Conditional on $S=\sum_iX_i=s$, symmetry gives
\[
    \Prob_p(X_1=1\mid S=s)=\frac{s}{n},
\]
which does not depend on $p$.  Hence $X_1$ can be simulated from $S$ by a
$\Bern(s/n)$ draw.  More importantly, replacing $X_1$ by its conditional mean
$S/n$ reduces squared-error risk.  The reduction $S$ therefore supports a
strictly better nonrandomized estimator than the original rule.
\end{example}

\section{Sufficiency}

\begin{definition}[Sufficient statistic]
A statistic $T=T(X)$ is \emph{sufficient} for $\theta$ (or for the family
$\mathcal P$) if there is a version of the conditional distribution of $X$
given $T$ that is the same for every $\theta\in\Theta$.  Equivalently, there
is a Markov kernel $K(\dd x\mid t)$, not depending on $\theta$, such that
$K(\cdot\mid T)$ is a version of $P_\theta(X\in\cdot\mid T)$ for every
$\theta$.
\end{definition}

For a discrete model this says that, whenever $P_\theta(T=t)>0$,
\begin{equation}
    P_\theta(X=x\mid T=t)
    \label{eq:conditional-sufficiency}
\end{equation}
does not depend on $\theta$.  The general definition is phrased in terms of
conditional distributions because conditioning on $T=t$ may mean
conditioning on an event of probability zero.

\begin{example}[The number of heads is sufficient]
For $X_1,\ldots,X_n\iid\Bern(p)$ and $S=\sum_iX_i$, every binary string $x$
with $s$ heads has probability $p^s(1-p)^{n-s}$.  Consequently,
\[
    P_p(X=x\mid S=s)
    =\frac{p^s(1-p)^{n-s}}{\binom ns p^s(1-p)^{n-s}}
    =\binom ns^{-1}
\]
for $\sum_i x_i=s$ (with the boundary cases interpreted on their supports).
Conditional on $S=s$, the locations of the heads are uniform over the
$\binom ns$ possibilities and contain no further information about $p$.
\end{example}

The conditional definition is intuitive, but direct calculation of a
conditional law can be cumbersome.  In dominated models the following
criterion is usually easier.

\begin{theorem}[Fisher--Neyman factorization theorem]
Suppose $\{P_\theta:\theta\in\Theta\}$ is dominated by a $\sigma$-finite
measure $\mu$, and write $f_\theta=\dd P_\theta/\dd\mu$.  Under the usual
regularity conditions on the sample spaces, $T$ is sufficient if and only if
there are nonnegative measurable functions $g_\theta$ and $h$ such that
\begin{equation}
    f_\theta(x)=g_\theta(T(x))h(x)
    \quad\text{for $\mu$-almost every }x,
    \label{eq:factorization-theorem}
\end{equation}
where $h$ does not depend on $\theta$.
\end{theorem}

\begin{proof}[Proof sketch]
If \eqref{eq:factorization-theorem} holds, then within a fiber
$\{x:T(x)=t\}$ the relative weights of possible observations are determined
by $h(x)$, independently of $\theta$.  Normalizing these weights gives a
parameter-free conditional law of $X$ given $T=t$.  Conversely, if this
conditional law is parameter-free, disintegrate $P_\theta$ into the marginal
law of $T$ and the common conditional law of $X$ given $T$; their densities
provide $g_\theta$ and $h$.  A fully general proof uses regular conditional
probabilities and the Radon--Nikodym theorem.
\end{proof}

\begin{example}[Bernoulli sample]
For $x\in\{0,1\}^n$,
\[
    f_p(x)=\prod_{i=1}^n p^{x_i}(1-p)^{1-x_i}
           =p^{\sum_i x_i}(1-p)^{n-\sum_i x_i}.
\]
This factors through $S=\sum_iX_i$, with $h(x)=1$, so $S$ is sufficient for
$p$.  Any one-to-one transformation of $S$ is also sufficient; for example,
the sample mean $\bar X=S/n$ is sufficient.
\end{example}

\begin{example}[Normal location and scale]
If $X_1,\ldots,X_n\iid\Normal(\mu,\sigma^2)$ with both parameters unknown,
then
\[
 f_{\mu,\sigma^2}(x)
 =(2\pi\sigma^2)^{-n/2}
 \exp\left\{-\frac{1}{2\sigma^2}
 \left(\sum_{i=1}^n x_i^2-2\mu\sum_{i=1}^n x_i+n\mu^2\right)\right\}.
\]
The factorization theorem shows that
$T(X)=(\sum_iX_i,\sum_iX_i^2)$ is sufficient for $(\mu,\sigma^2)$.  When
$n\geq2$, this is equivalent to retaining $(\bar X,\sum_i(X_i-\bar X)^2)$.
\end{example}

Sufficiency concerns the family of distributions, not the way in which the
parameter is named.  If $T$ is sufficient for $\theta$, then it is sufficient
for any function $g(\theta)$.  The converse need not hold if the model is
changed so that only $g(\theta)$ is regarded as unknown.

\subsection{Exponential Families}

A dominated family is a $k$-parameter \emph{exponential family} if its density
can be written
\begin{equation}
    f_\eta(x)
    =h(x)\exp\!\left\{\eta^\mathsf{T}T(x)-A(\eta)\right\},
    \qquad \eta\in H\subseteq\R^k.
    \label{eq:exponential-family}
\end{equation}
Here $\eta$ is the natural parameter, $T(x)\in\R^k$ is the natural statistic,
and
\[
    A(\eta)=\log\int h(x)e^{\eta^\mathsf{T}T(x)}\,\mu(\dd x)
\]
is the log-partition function.  The natural parameter space is
\[
    H_0=\left\{\eta\in\R^k:
    \int h(x)e^{\eta^\mathsf{T}T(x)}\,\mu(\dd x)<\infty\right\}.
\]
A family is called \emph{full} when its parameter space is $H_0$, and
\emph{regular} when $H$ is open.  Some authors use slightly different
terminology, so these conditions should always be stated explicitly.

For an iid sample with one-observation density
\[
    f_\eta(x)=h(x)\exp\{\eta^\mathsf{T}t(x)-A(\eta)\},
\]
the joint density is
\[
    \left\{\prod_{i=1}^nh(x_i)\right\}
    \exp\left\{\eta^\mathsf{T}\sum_{i=1}^nt(x_i)-nA(\eta)\right\}.
\]
Thus $\sum_i t(X_i)$ is sufficient, and its dimension does not grow with the
sample size.

\begin{example}[Bernoulli family in natural form]
For $0<p<1$, set $\eta=\log\{p/(1-p)\}$.  Then
\[
    p^x(1-p)^{1-x}
    =\exp\{x\eta-\log(1+e^\eta)\},\qquad x\in\{0,1\}.
\]
Hence $A(\eta)=\log(1+e^\eta)$ and the natural statistic for an iid sample is
$S=\sum_iX_i$.  The parameter space is all of $\R$, so this is a full,
regular, one-parameter exponential family.
\end{example}

\begin{example}[A curved exponential family]
Let $X_1,\ldots,X_n\iid\Normal(\mu,\mu^2)$ with $\mu>0$.  The joint density
factors through $(\sum_iX_i,\sum_iX_i^2)$, but its two natural parameters are
constrained by a one-dimensional relation.  It is therefore a curved
subfamily of a two-parameter exponential family.  Curved families need not
inherit all completeness properties of the containing full family.
\end{example}

The log-partition function encodes moments of the natural statistic.  If
differentiation under the integral is justified and $\eta$ is in the interior
of the natural parameter space, then
\begin{equation}
    \grad A(\eta)=\E_\eta[T(X)],
    \qquad
    \Hess A(\eta)=\Cov_\eta(T(X)).
    \label{eq:log-partition-moments}
\end{equation}
In particular $A$ is convex.  Strict convexity corresponds, after removing
redundant components, to a nondegenerate natural statistic.


\subsection{Minimal Sufficiency}

Sufficient statistics are not unique: if $T$ is sufficient, then
$(T,U)$ is sufficient for any statistic $U$.  We therefore seek the coarsest
sufficient reduction.

\begin{definition}[Minimal sufficient statistic]
A sufficient statistic $T$ is \emph{minimal sufficient} if, for every
sufficient statistic $S$, there exists a measurable function $q$ such that
\[
    T=q(S)\qquad P_\theta\text{-almost surely for every }\theta.
\]
Equivalently, the $\sigma$-field generated by $T$ is contained, up to common
null sets, in that generated by every other sufficient statistic.
\end{definition}

Minimality does not mean that $T$ has the smallest Euclidean dimension.  A
one-to-one encoding can map a vector into a scalar without changing its
information.  The relevant object is the partition of the sample space into
level sets of $T$, or equivalently the $\sigma$-field $\sigma(T)$.

\begin{theorem}[Likelihood-ratio criterion]
Suppose the model is dominated and its densities have common support
$\mathcal X_0$.  If a statistic $T$ satisfies
\begin{equation}
    T(x)=T(y)
    \quad\Longleftrightarrow\quad
    \frac{f_\theta(x)}{f_\theta(y)}
    \text{ is constant as a function of $\theta$},
    \label{eq:minimal-ratio-criterion}
\end{equation}
for $x,y\in\mathcal X_0$, then $T$ is minimal sufficient.  With suitable
measure-theoretic qualifications, the criterion remains valid beyond the
strict common-support setting.
\end{theorem}

\begin{proof}[Idea of proof]
The forward implication in \eqref{eq:minimal-ratio-criterion} yields a
factorization and hence sufficiency.  If $S$ is any sufficient statistic,
then $S(x)=S(y)$ forces the likelihood ratio $f_\theta(x)/f_\theta(y)$ to be
free of $\theta$.  The reverse implication in
\eqref{eq:minimal-ratio-criterion} then forces $T(x)=T(y)$.  Thus every level
set of $S$ lies inside a level set of $T$, so $T$ is a function of $S$.
\end{proof}

\begin{example}[Minimal sufficiency for coin flips]
For two strings $x,y\in\{0,1\}^n$ and $0<p<1$,
\[
    \frac{f_p(x)}{f_p(y)}
    =\left(\frac{p}{1-p}\right)^{\sum_i x_i-\sum_i y_i}.
\]
This ratio is constant in $p$ if and only if
$\sum_i x_i=\sum_i y_i$.  Therefore the number of heads is minimal
sufficient for $p$.
\end{example}

Parameter-dependent support requires care but often makes the criterion
especially informative.

\begin{example}[Uniform endpoint]
Let $X_1,\ldots,X_n\iid\mathrm{Uniform}(0,\theta)$, $\theta>0$.  The joint
density is
\[
    f_\theta(x)=\theta^{-n}
    \ind\{0\leq x_{(1)},\ x_{(n)}\leq\theta\}.
\]
Because the lower endpoint is known to be zero, this factors through the
sample maximum $X_{(n)}$.  For samples $x,y$ in the nonnegative orthant, the
likelihood ratio is independent of $\theta$ exactly when
$x_{(n)}=y_{(n)}$; if the maxima differ, there are values of $\theta$ for
which one likelihood is zero and the other is not.  Hence $X_{(n)}$ is
minimal sufficient.
\end{example}

\subsection{Complete Sufficiency}

Sufficiency says that a statistic retains all information in the sample.
Completeness says that its distribution contains enough variation across
$\theta$ to distinguish integrable functions of that statistic.

\begin{definition}[Completeness]
A statistic $T$ is \emph{complete} for $\mathcal P$ if, for every measurable
function $a$ for which $\E_\theta\abs{a(T)}<\infty$ for all $\theta$,
\[
    \E_\theta[a(T)]=0\quad\text{for every }\theta\in\Theta
    \qquad\Longrightarrow\qquad
    a(T)=0\quad P_\theta\text{-almost surely for every }\theta.
\]
It is \emph{boundedly complete} if the implication is required only for
bounded $a$.
\end{definition}

Completeness is a property of a family of distributions of $T$, not of a
single distribution.  It will later guarantee uniqueness of unbiased
estimators that are functions of a sufficient statistic.

\begin{theorem}[Completeness of full exponential families]
Consider the exponential family \eqref{eq:exponential-family}.  If its natural
parameter space contains a nonempty open subset of $\R^k$, then the natural
statistic $T$ is complete, subject to the stated integrability condition.
\end{theorem}

\begin{proof}[Proof sketch]
Suppose $\E_\eta[a(T)]=0$ throughout an open set.  After multiplying by
$e^{A(\eta)}$, this identity says that the Laplace transform of a signed
measure induced by $a(T)h(X)\mu(\dd x)$ vanishes on an open set.  Uniqueness
of Laplace transforms implies that this signed measure is zero, and therefore
$a(T)=0$ almost surely throughout the family.
\end{proof}

The openness condition matters.  If a $k$-dimensional exponential family is
restricted to a lower-dimensional curve, a nonzero function may have mean
zero at every parameter value in the restricted family.

\begin{example}[Completeness of the binomial count]
Let $S\sim\Bin(n,p)$, $0<p<1$, and suppose $\E_p[a(S)]=0$ for every $p$.
Then
\[
    0=\sum_{s=0}^n a(s)\binom ns p^s(1-p)^{n-s}.
\]
Dividing by $(1-p)^n$ and putting $z=p/(1-p)>0$ gives
\[
    \sum_{s=0}^n a(s)\binom ns z^s=0
    \qquad\text{for every }z>0.
\]
A polynomial that vanishes on an interval has every coefficient equal to
zero, so $a(s)=0$ for $s=0,\ldots,n$.  Thus $S$ is complete.  Together with
the preceding results, $S$ is complete and sufficient for $p$.
\end{example}

\begin{example}[Sufficient but not complete]
Let $X_1,\ldots,X_n\iid\mathrm{Uniform}(\theta-1,\theta+1)$.  The pair
$T=(X_{(1)},X_{(n)})$ is sufficient for $\theta$.  The range
$R=X_{(n)}-X_{(1)}$ is a nonconstant function of $T$, but its distribution
does not depend on $\theta$.  Hence
\[
    a(T)=R-\E_\theta[R]
\]
has expectation zero for every $\theta$ without being almost surely zero.
Therefore $T$ is not complete.
\end{example}

\begin{proposition}
If $T$ is complete and sufficient, then $T$ is minimal sufficient (up to
common null sets).
\end{proposition}

\begin{proof}[Proof sketch]
Let $S$ be any sufficient statistic and let $B\in\sigma(T)$.  By sufficiency
of $S$, a version of $q(S)=\E_\theta[\ind_B\mid S]$ can be chosen independent
of $\theta$.  By sufficiency of $T$ and iterated conditioning,
$\E_\theta[q(S)\mid T]$ is also parameter-free.  Its difference from
$\ind_B$ has mean zero for every $\theta$; completeness forces that
difference to vanish.  This ultimately shows that $B$ is measurable with
respect to $S$ modulo common null sets.  Hence $\sigma(T)\subseteq\sigma(S)$.
\end{proof}

\subsection{Rao-Blackwell Theorem}

Sufficiency is not only a statement about compression; it gives a systematic
way to improve statistical procedures.  Given an estimator $U=U(X)$ and a
sufficient statistic $T$, define its Rao--Blackwellization by
\[
    U^*(T)=\E_\theta[U\mid T].
\]
Because $T$ is sufficient, the conditional distribution of $X$ given $T$ is
parameter-free.  Consequently the right-hand side can be chosen as a single
function of $T$ that does not involve the unknown $\theta$.

\begin{theorem}[Rao--Blackwell]
Let $T$ be sufficient for $\theta$, let $U$ be an estimator of $g(\theta)$,
and suppose the relevant expectations exist.  If
$a\mapsto L(\theta,a)$ is convex for every $\theta$, then
\begin{equation}
    R(\theta,U^*)\leq R(\theta,U)
    \qquad\text{for every }\theta.
    \label{eq:rao-blackwell-risk}
\end{equation}
If the loss is strictly convex, equality at a given $\theta$ can hold only
when $U=U^*$ almost surely under $P_\theta$ (apart from degeneracies in the
loss).
\end{theorem}

\begin{proof}
Conditional Jensen's inequality gives
\[
 L\!\left(\theta,\E_\theta[U\mid T]\right)
 \leq \E_\theta[L(\theta,U)\mid T].
\]
Taking $P_\theta$-expectations proves \eqref{eq:rao-blackwell-risk}.
Sufficiency ensures that $U^*$ is a statistic rather than a rule involving
the unknown parameter.  The equality statement follows from the equality
condition in strict Jensen's inequality.
\end{proof}

Under squared-error loss, the result has an exact variance decomposition.  If
$U$ has finite second moment, then
\begin{equation}
    \Var_\theta(U)
    =\Var_\theta\!\left(\E_\theta[U\mid T]\right)
     +\E_\theta\!\left[\Var_\theta(U\mid T)\right].
    \label{eq:total-variance-rb}
\end{equation}
Moreover $\E_\theta[U^*]=\E_\theta[U]$, so Rao--Blackwellization preserves
bias.  Its mean squared error can only decrease, and the reduction equals
$\E_\theta[\Var_\theta(U\mid T)]$.

\begin{example}[One flip versus all flips]
Let $X_1,\ldots,X_n\iid\Bern(p)$.  The estimator $U=X_1$ is unbiased for $p$,
and $S=\sum_iX_i$ is sufficient.  Given $S=s$, each of the $n$ positions is
equally likely to contain any one of the $s$ heads, so
\[
    U^*=\E_p[X_1\mid S]=\frac{S}{n}=\bar X.
\]
The variances are
\[
    \Var_p(X_1)=p(1-p),
    \qquad
    \Var_p(\bar X)=\frac{p(1-p)}{n}.
\]
Thus using the sufficient statistic reduces variance by a factor of $n$.
\end{example}

\begin{example}[Estimating the probability of all heads]
In the same model, $U=\prod_{i=1}^nX_i$ is unbiased for $p^n$.  Since $U=1$
exactly when $S=n$, it is already a function of the sufficient statistic and
Rao--Blackwellization makes no change.  By contrast, for $m\leq n$,
$U_m=\prod_{i=1}^mX_i$ is unbiased for $p^m$ and
\[
    \E_p[U_m\mid S=s]
    =\frac{\binom{s}{m}}{\binom{n}{m}},
\]
where the ratio is defined as zero for $s<m$.  This is the conditional
probability that a fixed set of $m$ positions is contained among the $s$
heads, and it has no dependence on $p$.
\end{example}

Rao--Blackwellization need not produce a unique or globally optimal estimator
by itself.  Completeness supplies the missing uniqueness.

\begin{theorem}[Lehmann--Scheff\'e]
If $T$ is complete and sufficient and $U$ is any unbiased estimator of
$g(\theta)$ with finite variance, then
\[
    \delta(T)=\E_\theta[U\mid T]
\]
is the unique unbiased estimator of $g(\theta)$ that is a function of $T$,
and it has variance no larger than that of any unbiased estimator.  Thus it is
the uniformly minimum variance unbiased estimator of $g(\theta)$.
\end{theorem}

\begin{proof}
Rao--Blackwellization shows that $\delta(T)$ is unbiased and has no larger
variance than $U$.  If $\delta_1(T)$ and $\delta_2(T)$ are two unbiased
functions of $T$, then their difference has expectation zero for all
$\theta$.  Completeness implies $\delta_1(T)=\delta_2(T)$ almost surely.
Applying Rao--Blackwellization to any competing unbiased estimator proves the
minimum-variance claim.
\end{proof}

For the coin model, the binomial count $S$ is complete and sufficient.
Consequently $S/n$ is the unique uniformly minimum variance unbiased
estimator of $p$, and
$\binom{S}{m}/\binom{n}{m}$ is the corresponding estimator of $p^m$.  These
examples summarize the main data-reduction pipeline:
\[
    \text{raw sample}
    \longrightarrow \text{sufficient statistic}
    \longrightarrow \text{Rao--Blackwellized rule},
\]
with completeness converting the improved rule into a uniqueness and
optimality statement.

\part{Optimal Estimation, Exact}

\section{Uniformly Minimum Variance Unbiased Estimators}

Let $g:\Theta\to\R$ be the estimand.  An estimator $\delta(X)$ is
\emph{unbiased} for $g(\theta)$ if
\[
    \E_\theta[\delta(X)]=g(\theta)
    \qquad\text{for every }\theta\in\Theta.
\]
Under squared-error loss, its risk decomposes as
\begin{equation}
    R(\theta,\delta)
    =\E_\theta[(\delta-g(\theta))^2]
    =\Var_\theta(\delta)
      +\{\E_\theta[\delta]-g(\theta)\}^2.
    \label{eq:bias-variance-decomposition}
\end{equation}
Thus risk equals variance within the class of unbiased estimators.

\begin{definition}[UMVU estimator]
An unbiased estimator $\delta_0$ of $g(\theta)$ is a \emph{uniformly minimum
variance unbiased} (UMVU) estimator if
\[
    \Var_\theta(\delta_0)
    \leq \Var_\theta(\delta)
\]
for every $\theta$ and every unbiased estimator $\delta$ of $g(\theta)$.
The word ``uniformly'' means that the same estimator minimizes variance at
every parameter value.
\end{definition}

Pointwise minimization is not automatic: an estimator that has smaller
variance near one parameter value may have larger variance elsewhere.  A
useful characterization turns the minimization problem into an orthogonality
condition.

\begin{theorem}[Orthogonality characterization]
Suppose $\delta_0$ is unbiased for $g(\theta)$ and has finite second moment.
Then $\delta_0$ is UMVU if and only if
\begin{equation}
    \Cov_\theta(\delta_0,U)=0
    \label{eq:umvu-orthogonality}
\end{equation}
for every finite-variance estimator $U$ satisfying
$\E_\theta[U]=0$ for all $\theta$.
\end{theorem}

\begin{proof}
For every real $a$, $\delta_0+aU$ is unbiased and
\[
 \Var_\theta(\delta_0+aU)
 =\Var_\theta(\delta_0)+2a\Cov_\theta(\delta_0,U)
  +a^2\Var_\theta(U).
\]
If $\delta_0$ is UMVU, this quadratic is minimized at $a=0$, forcing the
covariance to vanish.  Conversely, the difference $U=\delta-\delta_0$ between
any two unbiased estimators has mean zero.  Orthogonality then gives
$\Var_\theta(\delta)=\Var_\theta(\delta_0)+\Var_\theta(U)$.
\end{proof}

The main construction is the Lehmann--Scheff\'e theorem from the preceding
part: find a complete sufficient statistic $T$, begin with any unbiased
estimator, and condition on $T$.

\begin{example}[Powers of a coin bias]
Let $S\sim\Bin(n,p)$.  Since $S$ is complete and sufficient, the UMVU
estimator of $p$ is $S/n$.  More generally, for $1\leq m\leq n$,
\[
    \widehat{p^m}=\frac{(S)_m}{(n)_m}
    =\frac{\binom Sm}{\binom nm},
\]
where $(a)_m=a(a-1)\cdots(a-m+1)$.  Indeed,
$\E_p[(S)_m]=(n)_mp^m$.  No unbiased estimator of $p^m$ exists when $m>n$:
the expectation of any function of $S$ is a polynomial in $p$ of degree at
most $n$.
\end{example}

\begin{example}[Normal mean and variance]
If $X_1,\ldots,X_n\iid\Normal(\mu,\sigma^2)$ with both parameters unknown,
then $(\sum_iX_i,\sum_iX_i^2)$ is complete and sufficient because the normal
family is a full-rank exponential family with an open natural parameter
space.  Therefore $\bar X$ is UMVU for $\mu$, and
\[
    S_X^2=\frac{1}{n-1}\sum_{i=1}^n(X_i-\bar X)^2
\]
is UMVU for $\sigma^2$.
\end{example}

\section{Uniformly Minimum Risk Unbiased Estimators}

Variance is tied specifically to scalar estimation under squared-error loss.
For a general loss, an analogous restricted optimum is useful.

\begin{definition}[UMRU estimator]
Within a specified class $\mathcal U$ of unbiased decision rules, a rule
$\delta_0$ is \emph{uniformly minimum risk unbiased} (UMRU) if
\[
    R(\theta,\delta_0)\leq R(\theta,\delta)
    \qquad\text{for every }\theta\in\Theta,
    \quad \delta\in\mathcal U.
\]
For scalar estimation, $\mathcal U$ usually consists of rules satisfying
$\E_\theta[\delta]=g(\theta)$.
\end{definition}

Under squared-error loss, UMRU and UMVU are the same notion by
\eqref{eq:bias-variance-decomposition}.  Under another loss they can differ.
For example, unbiasedness constrains a mean, whereas absolute-error risk is
governed by the entire distribution and is naturally minimized by medians.

There is also an intrinsic definition of unbiasedness for a decision problem
that does not require a vector-valued action to have a mean.

\begin{definition}[Loss unbiasedness]
A decision rule $\delta$ is \emph{risk unbiased}, or \emph{unbiased in loss},
if
\begin{equation}
    \E_\theta[L(\theta,\delta(X))]
    \leq \E_\theta[L(\theta',\delta(X))]
    \quad\text{for all }\theta,\theta'\in\Theta.
    \label{eq:risk-unbiasedness}
\end{equation}
Thus, when the data come from $P_\theta$, the rule is on average at least as
close to the true target $\theta$ as to any false target $\theta'$.
\end{definition}

For $L(\theta,a)=(a-\theta)^2$ and finite second moments, subtracting the two
sides of \eqref{eq:risk-unbiasedness} shows that it holds for every
$\theta'$ if and only if $\E_\theta[\delta]=\theta$.  Hence ordinary and risk
unbiasedness agree in this important case.

If $T$ is sufficient and the loss is convex in the action, conditioning a
rule on $T$ preserves ordinary unbiasedness and reduces risk by conditional
Jensen.  Consequently the search for a UMRU rule may be restricted to
functions of $T$.  Completeness then provides uniqueness when the relevant
unbiasedness equations reduce to expectations of functions of $T$.

\section{Information Inequality}

The information inequality supplies a lower bound on the variance of any
unbiased estimator without explicitly comparing it with all competitors.
Assume first that $\theta$ is scalar and that $P_\theta$ has density
$f_\theta$ with respect to a fixed measure.  The \emph{score} and Fisher
information are
\[
    \ell'_\theta(X)=\frac{\partial}{\partial\theta}
        \log f_\theta(X),
    \qquad
    I_X(\theta)=\E_\theta[(\ell'_\theta(X))^2].
\]
Under regularity conditions allowing differentiation under the integral and
requiring parameter-independent support,
\[
    \E_\theta[\ell'_\theta(X)]=0,
    \qquad
    I_X(\theta)=-\E_\theta[\ell''_\theta(X)].
\]

\begin{theorem}[Cram\'er--Rao information inequality]
Suppose $\delta$ has finite variance,
$\E_\theta[\delta]=g(\theta)$, and the usual differentiation and support
conditions hold.  If $0<I_X(\theta)<\infty$, then
\begin{equation}
    \Var_\theta(\delta)
    \geq \frac{\{g'(\theta)\}^2}{I_X(\theta)}.
    \label{eq:cramer-rao-scalar}
\end{equation}
Equality holds at $\theta$ precisely when
$\delta-g(\theta)$ is almost surely proportional to the score.
\end{theorem}

\begin{proof}
Differentiating $\E_\theta[\delta]=g(\theta)$ gives
\[
 g'(\theta)=\E_\theta[\delta\ell'_\theta]
 =\Cov_\theta(\delta,\ell'_\theta).
\]
Cauchy--Schwarz implies
$\{g'(\theta)\}^2\leq\Var_\theta(\delta)I_X(\theta)$.  Its equality condition
gives the final assertion.
\end{proof}

For iid observations, scores add and have mean zero, so information adds:
$I_{X_1,\ldots,X_n}(\theta)=nI_{X_1}(\theta)$.

\begin{example}[Weighted coin]
For one $\Bern(p)$ observation,
\[
    \ell'_p(X)=\frac{X-p}{p(1-p)},
    \qquad I_X(p)=\frac{1}{p(1-p)}.
\]
For $n$ flips, every unbiased estimator of $p$ has variance at least
$p(1-p)/n$.  The sample proportion $S/n$ attains the bound because
\[
    \frac Sn-p=\frac{p(1-p)}{n}\ell'_p(X_1,\ldots,X_n).
\]
It is therefore efficient and UMVU.
\end{example}

If $\theta\in\R^k$, define the score vector and information matrix by
\[
    s_\theta(X)=\grad_\theta\log f_\theta(X),
    \qquad I_X(\theta)=\E_\theta[s_\theta s_\theta^\mathsf T].
\]
For an unbiased estimator of a scalar $g(\theta)$, the matrix form is
\begin{equation}
    \Var_\theta(\delta)
    \geq \grad g(\theta)^\mathsf T I_X(\theta)^{-1}
       \grad g(\theta).
    \label{eq:cramer-rao-matrix}
\end{equation}
For a vector estimator $D$ with mean $g(\theta)$, the corresponding covariance
matrix inequality is
\[
 \Cov_\theta(D)\succeq
 \mathrm Dg(\theta)I_X(\theta)^{-1}\mathrm Dg(\theta)^\mathsf T,
\]
where $A\succeq B$ means that $A-B$ is positive semidefinite.

The regularity assumptions cannot be ignored.  For
$X_1,\ldots,X_n\iid\mathrm{Uniform}(0,\theta)$, the support depends on
$\theta$, and differentiating the integral as above misses a boundary term.
Indeed, a rescaled sample maximum is unbiased and can have variance smaller
than a naive application of the regular Cram\'er--Rao formula would permit.

\section{Minimum Risk Equivariant Estimators}

Many models possess symmetries.  Let a group $G$ act on the sample space,
parameter space, and action space.  A problem is \emph{invariant} if
transforming both the data and parameter leaves the model and loss unchanged:
\[
    X\sim P_\theta \Longrightarrow gX\sim P_{g\theta},
    \qquad L(g\theta,ga)=L(\theta,a).
\]
A rule $\delta$ is \emph{equivariant} if
\[
    \delta(gx)=g\delta(x).
\]
For an equivariant rule, risk is constant along group orbits:
$R(g\theta,\delta)=R(\theta,\delta)$.  When the group acts transitively on
$\Theta$, its risk is constant over the whole parameter space.  A
\emph{minimum risk equivariant} (MRE) rule minimizes this constant among all
equivariant rules.

\subsection{Location and Scale Families}

In a location family, $X_i=\mu+Z_i$ for a known distribution of the errors.
The translation group sends $x$ to $x+c\bm 1$ and $\mu$ to $\mu+c$.  Under
loss $L(\mu,a)=\rho(a-\mu)$, an equivariant estimator satisfies
\[
    \delta(x+c\bm 1)=\delta(x)+c.
\]
Its error distribution does not depend on $\mu$, and hence its risk is
constant.

Under squared-error loss, the MRE location estimator is the generalized Bayes
rule under the translation-invariant prior $\pi(\mu)\propto1$:
\begin{equation}
    \delta_{\mathrm{MRE}}(x)
    =\frac{\int_{-\infty}^{\infty}\mu f(x-\mu\bm 1)\dd\mu}
           {\int_{-\infty}^{\infty}f(x-\mu\bm 1)\dd\mu},
    \label{eq:pitman-estimator}
\end{equation}
provided these integrals exist.  This is often called the Pitman estimator.
For normal errors with known variance it reduces to $\bar X$.

In a positive scale family, $X_i=\sigma Z_i$, $\sigma>0$.  The scale group
sends $x$ to $cx$, $\sigma$ to $c\sigma$, and an equivariant estimator of
$\sigma$ obeys $\delta(cx)=c\delta(x)$.  A scale-invariant loss, such as
\[
    L(\sigma,a)=\left(\frac a\sigma-1\right)^2,
\]
makes the risk of every scale-equivariant rule constant.  The invariant
measure on the scale group leads to the improper prior
$\pi(\sigma)\propto1/\sigma$; the resulting generalized Bayes rule is a
candidate MRE estimator.

Improper priors do not define probability distributions, so generalized
Bayes calculations require separate justification.  The posterior integrals
must be finite, and optimality among equivariant rules does not by itself
imply optimality among all rules.

\subsection{Risk Unbiasedness}

Equivariance and unbiasedness impose different restrictions.  Equivariance
requires a rule to respect the geometry of the problem; risk unbiasedness
requires the true parameter to minimize expected loss as in
\eqref{eq:risk-unbiasedness}.  Neither condition generally implies the other.

For a location problem under squared-error loss, every equivariant rule can
be written schematically as
\[
    \delta(X)=X_1+h(X_2-X_1,\ldots,X_n-X_1).
\]
Its bias is constant in $\mu$.  It is risk unbiased exactly when that constant
is zero.  Thus risk unbiasedness may select a proper subset of the equivariant
class, while MRE optimality selects the rule with smallest constant risk in
the larger class.

The distinction matters in constrained parameter spaces.  If a parameter is
known to lie in an interval, an estimator can often be improved under squared
error by projecting it onto that interval.  Projection generally introduces
bias but decreases loss pointwise.  Hence restricting attention to unbiased
rules can exclude sensible procedures.

\section{Average Risk Optimality}

A risk function gives one value for every $\theta$, and two risk functions may
cross.  One way to obtain a scalar criterion is to average risk with respect
to a probability distribution $\pi$ on $\Theta$.  The \emph{Bayes risk} is
\begin{equation}
    r(\pi,\delta)=\int_\Theta R(\theta,\delta)\,\pi(\dd\theta).
    \label{eq:bayes-risk}
\end{equation}
This does not require a subjective interpretation: $\pi$ can also be viewed
as a device for weighting regions of the parameter space.

\subsection{Bayes Estimators}

By Tonelli's theorem, the Bayes risk can be written using the posterior:
\[
 r(\pi,\delta)
 =\int_{\mathcal X}
   \left\{\int_\Theta L(\theta,\delta(x))\pi(\dd\theta\mid x)\right\}
   m_\pi(\dd x),
\]
where $m_\pi$ is the prior predictive distribution.  Therefore a Bayes rule
may be found by minimizing posterior expected loss separately for each $x$.

\begin{proposition}[Bayes actions]
For a scalar target $g(\theta)$, a Bayes estimator is
\begin{itemize}
    \item the posterior mean under squared-error loss;
    \item any posterior median under absolute-error loss;
    \item a posterior mode under zero-one loss in a discrete action space.
\end{itemize}
\end{proposition}

\begin{example}[Beta--binomial estimation]
Let $S\mid p\sim\Bin(n,p)$ and $p\sim\BetaDist(\alpha,\beta)$, where
$\alpha,\beta>0$.  Conjugacy gives
\[
    p\mid S=s\sim\BetaDist(\alpha+s,\beta+n-s).
\]
Under squared-error loss the Bayes estimator is
\begin{equation}
    \delta_\pi(s)=\E[p\mid S=s]
    =\frac{\alpha+s}{\alpha+\beta+n}
    =w\frac{s}{n}+(1-w)\frac{\alpha}{\alpha+\beta},
    \label{eq:beta-binomial-bayes}
\end{equation}
where $w=n/(n+\alpha+\beta)$.  It shrinks the observed proportion toward the
prior mean.  Its posterior squared-error loss is minimized uniquely, and its
Bayes risk equals the expected posterior variance.
\end{example}

\begin{theorem}[Admissibility of a unique Bayes rule]
If a Bayes rule $\delta_\pi$ has finite Bayes risk, is unique up to prior
predictive null sets, and every nonempty set on which a competing rule could
strictly improve has positive prior mass, then $\delta_\pi$ is admissible.
For example, the last condition follows when $\pi$ has full support and the
risk difference is continuous in $\theta$.
\end{theorem}

\begin{proof}
If a rule dominated $\delta_\pi$, integrating the risk inequality against
$\pi$ would give no larger Bayes risk.  Strict improvement on a set of
positive prior mass would give strictly smaller Bayes risk, contradicting
Bayes optimality.  Uniqueness rules out domination that differs only on sets
invisible to the prior predictive distribution.
\end{proof}

\section{Worst Case Optimality}

When no averaging distribution is accepted, a conservative criterion is the
maximum risk
\[
    \overline R(\delta)=\sup_{\theta\in\Theta}R(\theta,\delta).
\]
This treats nature as selecting the least favorable parameter after the rule
has been chosen.

\subsection{Minimax Estimators}

\begin{definition}[Minimax rule]
A rule $\delta^*$ is minimax if
\[
    \sup_\theta R(\theta,\delta^*)
    =\inf_\delta\sup_\theta R(\theta,\delta).
\]
The common value is the minimax value of the decision problem.
\end{definition}

For every prior $\pi$ and rule $\delta$,
\begin{equation}
    \inf_{\widetilde\delta}\sup_\theta R(\theta,\widetilde\delta)
    \geq \inf_{\widetilde\delta}r(\pi,\widetilde\delta)
    =r(\pi,\delta_\pi).
    \label{eq:bayes-lower-minimax}
\end{equation}
Thus every Bayes risk is a lower bound on the minimax value.

\begin{theorem}[Constant-risk Bayes criterion]
If a Bayes rule $\delta_\pi$ has constant risk $R(\theta,\delta_\pi)=c$ for
all $\theta$ in the parameter space, then it is minimax and its Bayes risk is
$c$.
\end{theorem}

\begin{proof}
The maximum risk of $\delta_\pi$ is $c$.  On the other hand,
\eqref{eq:bayes-lower-minimax} says that the minimax value is at least its
Bayes risk, also $c$.
\end{proof}

The theorem explains why invariant rules with constant risk are natural
minimax candidates.  A generalized Bayes rule from an improper invariant
prior cannot be inserted directly into the proof, but it is often obtained as
a limit of proper Bayes rules whose Bayes risks approach its constant risk.

\subsection{Least Favorable Priors}

\begin{definition}[Least favorable prior]
A prior $\pi^*$ is \emph{least favorable} if it maximizes the optimal Bayes
risk:
\[
    r(\pi^*,\delta_{\pi^*})
    =\sup_\pi\inf_\delta r(\pi,\delta).
\]
It concentrates weight where the estimation problem is hardest.
\end{definition}

If a prior $\pi^*$ and its Bayes rule satisfy
\[
    \sup_\theta R(\theta,\delta_{\pi^*})
    =r(\pi^*,\delta_{\pi^*}),
\]
then the lower and upper bounds meet: $\delta_{\pi^*}$ is minimax and
$\pi^*$ is least favorable.  In finite games this is the familiar saddle
point condition.

Often no proper least favorable prior exists.  It is enough to find priors
$\pi_m$ whose Bayes risks converge to the maximum risk of a candidate rule:
if
\[
    r(\pi_m,\delta_{\pi_m})\longrightarrow c
    \quad\text{and}\quad
    \sup_\theta R(\theta,\delta^*)=c,
\]
then \eqref{eq:bayes-lower-minimax} proves that $\delta^*$ is minimax.  Such a
sequence is called a least favorable sequence.

\begin{example}[Normal location]
If $X\sim\Normal(\theta,\sigma^2)$ and loss is squared error, the estimator
$\delta(X)=X$ has constant risk $\sigma^2$.  Take proper priors
$\theta\sim\Normal(0,\tau^2)$.  The Bayes rule is
\[
    \delta_\tau(X)=\frac{\tau^2}{\tau^2+\sigma^2}X,
\]
and its Bayes risk is
$\sigma^2\tau^2/(\sigma^2+\tau^2)\uparrow\sigma^2$ as $\tau^2\to\infty$.
Hence these priors form a least favorable sequence and $X$ is minimax.
\end{example}

\section{Simultaneous Estimation}

Suppose $X=(X_1,\ldots,X_d)^\mathsf T\sim\Normal_d(\theta,I_d)$ and the goal is
to estimate the entire mean vector under total squared-error loss
\[
    L(\theta,a)=\norm{a-\theta}^2.
\]
The coordinatewise estimator $\delta_0(X)=X$ is unbiased, equivariant, and
has constant risk
\[
    R(\theta,\delta_0)=\E_\theta\norm{X-\theta}^2=d.
\]
It is minimax, since it is the limiting Bayes rule under increasingly diffuse
normal priors.  Surprisingly, minimaxity does not imply admissibility.

\subsection{James-Stein Estimator}

For $d\geq3$, define
\begin{equation}
    \delta_{\mathrm{JS}}(X)
    =\left(1-\frac{d-2}{\norm X^2}\right)X.
    \label{eq:james-stein}
\end{equation}
This estimator shrinks $X$ toward the origin.  Although every coordinate of
$X$ is individually the usual estimator of its mean, estimating the
coordinates jointly permits an improvement in total risk.

\begin{lemma}[Stein's identity]
If $X\sim\Normal_d(\theta,I_d)$ and $g:\R^d\to\R^d$ is sufficiently regular
and integrable, then
\[
    \E_\theta[(X-\theta)^\mathsf Tg(X)]
    =\E_\theta[\operatorname{div}g(X)].
\]
\end{lemma}

\begin{proof}
Apply one-dimensional integration by parts to each coordinate and sum.  The
normal density satisfies
$\partial f_\theta(x)/\partial x_i=-(x_i-\theta_i)f_\theta(x)$, and the
boundary terms vanish under the integrability assumptions.
\end{proof}

For an estimator $X+g(X)$, Stein's identity gives the unbiased risk formula
\begin{equation}
 R(\theta,X+g)-d
 =\E_\theta\left[\norm{g(X)}^2+2\operatorname{div}g(X)\right].
 \label{eq:stein-risk-identity}
\end{equation}
Set $g(x)=-ax/\norm x^2$.  Away from the origin,
\[
    \norm{g(x)}^2=\frac{a^2}{\norm x^2},
    \qquad
    \operatorname{div}g(x)=-\frac{a(d-2)}{\norm x^2}.
\]
Therefore
\[
 R(\theta,\delta_a)-d
 =\{a^2-2a(d-2)\}
   \E_\theta\left[\frac{1}{\norm X^2}\right].
\]
For $d\geq3$ this expectation is finite, and every
$0<a<2(d-2)$ strictly improves on $X$ for all $\theta$.  The choice $a=d-2$
gives \eqref{eq:james-stein} and maximizes the displayed risk reduction.

The positive-part version
\[
    \delta_{\mathrm{JS}}^+(X)
    =\left(1-\frac{d-2}{\norm X^2}\right)_+X
\]
dominates the ordinary James--Stein rule.  Shrinkage toward another fixed
point $m$ follows by replacing $X$ with $X-m$ and adding $m$ back.  The choice
of center expresses structural information and can have a substantial effect
on realized error, even though the dominance theorem holds for every fixed
center.

This phenomenon does not contradict the scalar Cram\'er--Rao bound: the
James--Stein estimator is biased, and performance is measured by the sum of
squared errors across coordinates.  It also shows why admissibility, UMVU
optimality, and minimaxity are distinct concepts.

\section{Empirical Bayes Estimators}

Hierarchical models provide a principled source of shrinkage.  Suppose
$\theta_1,\ldots,\theta_m$ are drawn from a distribution $G_\lambda$ with
unknown hyperparameter $\lambda$, and conditionally independent observations
satisfy $X_i\mid\theta_i\sim P_{\theta_i}$.  If $\lambda$ were known, the Bayes
rule would be
\[
    \delta_\lambda(x_i)=\E_\lambda[\theta_i\mid X_i=x_i]
\]
under squared-error loss.  An \emph{empirical Bayes} rule estimates
$\lambda$ from the ensemble $X_1,\ldots,X_m$ and plugs the estimate into the
Bayes rule.

\begin{example}[Many weighted coins]
For coin $i$, observe $S_i\mid p_i\sim\Bin(n_i,p_i)$ and assume
$p_i\iid\BetaDist(\alpha,\beta)$.  If $\alpha$ and $\beta$ were known, then
\[
    \E[p_i\mid S_i]
    =\frac{\alpha+S_i}{\alpha+\beta+n_i}.
\]
Estimate $(\alpha,\beta)$ from the marginal beta-binomial likelihood of all
counts, or by matching the across-coin mean and variance.  The empirical
Bayes estimator is
\[
    \widehat p_i^{\mathrm{EB}}
    =\frac{\widehat\alpha+S_i}
           {\widehat\alpha+\widehat\beta+n_i}.
\]
Coins with few flips are shrunk strongly toward the estimated population mean
$\widehat\alpha/(\widehat\alpha+\widehat\beta)$; coins with many flips are
driven mainly by their own sample proportions.
\end{example}

\begin{example}[Normal means]
Suppose $X_i\mid\theta_i\sim\Normal(\theta_i,\sigma^2)$ and
$\theta_i\iid\Normal(\mu,\tau^2)$, with $\sigma^2$ known.  The oracle Bayes
estimator is
\[
 \E[\theta_i\mid X_i]
 =\mu+B(X_i-\mu),
 \qquad B=\frac{\tau^2}{\tau^2+\sigma^2}.
\]
Replacing $\mu$ and $\tau^2$ by estimates from the marginal model
$X_i\sim\Normal(\mu,\tau^2+\sigma^2)$ yields an empirical Bayes shrinkage
rule.  A simple moment estimate is
\[
    \widehat\mu=\bar X,
    \qquad
    \widehat\tau^2=\max\left\{0,
       \frac{1}{m-1}\sum_{i=1}^m(X_i-\bar X)^2-\sigma^2\right\}.
\]
\end{example}

Empirical Bayes procedures borrow strength across related problems, but their
analysis must account for estimation of the prior.  Plug-in posterior
variances generally understate uncertainty because they treat
$\widehat\lambda$ as fixed.  Parametric bootstrap, asymptotic corrections, or
a fully Bayesian hyperprior can propagate this additional uncertainty.
Moreover, shrinkage is only as appropriate as the exchangeability assumption:
if the units come from substantially different populations, pooling can
introduce systematic bias.  These considerations separate the construction
of an empirical Bayes point estimator from the calibration of its uncertainty
statements.

\part{Optimal Hypothesis Testing, Exact}

\section{Introduction to Hypothesis Testing}

A hypothesis divides the parameter space into a null set $\Theta_0$ and an
alternative set $\Theta_1$.  We write
\[
    H_0:\theta\in\Theta_0
    \qquad\text{versus}\qquad
    H_1:\theta\in\Theta_1.
\]
A nonrandomized test is a measurable map $\phi:\mathcal X\to\{0,1\}$, where
$\phi(X)=1$ means that $H_0$ is rejected.  It is convenient to allow
randomized tests $\phi:\mathcal X\to[0,1]$, interpreted as the conditional
probability of rejection after observing $X=x$.

The \emph{power function} of a test is
\begin{equation}
    \beta_\phi(\theta)=\E_\theta[\phi(X)].
    \label{eq:power-function}
\end{equation}
For $\theta\in\Theta_0$, this is the probability of a type I error; for
$\theta\in\Theta_1$, $1-\beta_\phi(\theta)$ is the probability of a type II
error.  The \emph{size} is
\[
    \sup_{\theta\in\Theta_0}\beta_\phi(\theta),
\]
and a test has \emph{level $\alpha$} if its size is at most $\alpha$.  Size
equals level only when the supremum is attained at $\alpha$.

Testing is a decision problem with actions $0$ and $1$.  If false rejection
has loss $c_0$ and false acceptance has loss $c_1$, then
\[
 R(\theta,\phi)=
 \begin{cases}
   c_0\beta_\phi(\theta),&\theta\in\Theta_0,\\
   c_1\{1-\beta_\phi(\theta)\},&\theta\in\Theta_1.
 \end{cases}
\]
The frequentist level-$\alpha$ formulation fixes an upper bound on the first
kind of risk and then seeks large power on the alternative.

\begin{example}[Testing whether a coin is fair]
For $X_1,\ldots,X_n\iid\Bern(p)$, consider
$H_0:p=1/2$ against $H_1:p>1/2$.  Large values of
$S=\sum_iX_i$ constitute evidence against $H_0$.  A test of the form
\[
    \phi(s)=\ind\{s>c\}+\gamma\ind\{s=c\}
\]
can attain any desired level $\alpha$: choose $c$ so that
$P_{1/2}(S>c)\leq\alpha\leq P_{1/2}(S\geq c)$ and set
\[
    \gamma=\frac{\alpha-P_{1/2}(S>c)}{P_{1/2}(S=c)}.
\]
Randomization is needed only because the binomial distribution is discrete.
In practice, one often uses the conservative nonrandomized test with
$\gamma=0$.
\end{example}

A $p$-value is a statistic measuring how extreme the observation is relative
to a null distribution.  For a test statistic $T$ where large values are
evidence against $H_0$, a valid $p$-value $P(X)$ satisfies
\[
    \sup_{\theta\in\Theta_0}P_\theta\{P(X)\leq u\}\leq u,
    \qquad 0\leq u\leq1.
\]
Rejecting when $P(X)\leq\alpha$ therefore gives a level-$\alpha$ test.  A
$p$-value is not the posterior probability that the null is true.

\section{Simple Hypotheses}

Both hypotheses are \emph{simple} when they specify individual distributions:
\[
    H_0:\theta=\theta_0
    \qquad\text{versus}\qquad
    H_1:\theta=\theta_1.
\]
This case is completely solved by ordering observations according to their
likelihood ratio.

\subsection{Most Powerful (MP) Tests}

\begin{definition}[Most powerful test]
A level-$\alpha$ test $\phi^*$ is \emph{most powerful} for testing
$P_{\theta_0}$ against $P_{\theta_1}$ if
\[
    \E_{\theta_1}[\phi^*]
    \geq \E_{\theta_1}[\phi]
\]
for every other level-$\alpha$ test $\phi$.
\end{definition}

The level constraint is essential.  Without it, the rule that always rejects
has power one but provides no control of false rejection.  Randomization also
makes the feasible class convex and permits exact size in discrete models.

\begin{example}[One weighted-coin flip]
Test $H_0:p=p_0$ against $H_1:p=p_1$, where $p_1>p_0$, based on a single
flip.  Heads has likelihood ratio $p_1/p_0$, whereas tails has ratio
$(1-p_1)/(1-p_0)$; the former is larger.  For $\alpha\leq p_0$, the MP test
rejects with probability $\alpha/p_0$ after heads and never after tails.  For
$\alpha>p_0$, it always rejects after heads and rejects after tails with
probability $(\alpha-p_0)/(1-p_0)$.
\end{example}

\subsection{Neyman-Pearson Lemma}

\begin{theorem}[Neyman--Pearson lemma]
Suppose $P_0$ and $P_1$ have densities $f_0$ and $f_1$ with respect to a
common dominating measure.  Choose $k\geq0$ and $\gamma\in[0,1]$ so that
\begin{equation}
 \phi^*(x)=
 \begin{cases}
  1,&f_1(x)>kf_0(x),\\
  \gamma,&f_1(x)=kf_0(x),\\
  0,&f_1(x)<kf_0(x)
 \end{cases}
 \quad\text{satisfies}\quad \E_0[\phi^*]=\alpha.
 \label{eq:np-test}
\end{equation}
Then $\phi^*$ is most powerful among all tests of level $\alpha$.  Conversely,
every MP test of size $\alpha$ has this form, apart from null sets and possible
freedom on the boundary, provided the two densities are handled on their
zero sets in the usual way.
\end{theorem}

\begin{proof}
For any level-$\alpha$ test $\phi$,
\[
 \int(\phi^*-\phi)(f_1-kf_0)\dd\mu\geq0,
\]
because both factors have the same sign pointwise.  Hence
\[
 \E_1[\phi^*]-\E_1[\phi]
 \geq k\{\E_0[\phi^*]-\E_0[\phi]\}\geq0.
\]
This proves sufficiency.  The equality conditions in these inequalities yield
the converse characterization.
\end{proof}

\begin{example}[Several coin flips]
For $S\sim\Bin(n,p)$ and $p_1>p_0$,
\[
 \frac{f_{p_1}(s)}{f_{p_0}(s)}
 =\left(\frac{1-p_1}{1-p_0}\right)^n
  \left\{\frac{p_1(1-p_0)}{p_0(1-p_1)}\right\}^{s},
\]
which is strictly increasing in $s$.  Therefore an MP test rejects for large
numbers of heads, with possible randomization at one boundary count.
\end{example}

\subsection{Likelihood Ratio Tests}

For a general null and alternative, define the generalized likelihood ratio
statistic
\begin{equation}
    \Lambda(x)=
    \frac{\sup_{\theta\in\Theta_0}f_\theta(x)}
         {\sup_{\theta\in\Theta}f_\theta(x)},
    \qquad 0\leq\Lambda(x)\leq1.
    \label{eq:generalized-lr}
\end{equation}
The \emph{likelihood ratio test} (LRT) rejects for small $\Lambda$.  For two
simple hypotheses this is equivalent to the Neyman--Pearson likelihood-ratio
ordering.  For composite hypotheses, however, the Neyman--Pearson lemma does
not automatically make the LRT optimal.

\begin{example}[Two-sided normal mean]
If $X_1,\ldots,X_n\iid\Normal(\mu,\sigma^2)$ with $\sigma^2$ known, consider
$H_0:\mu=\mu_0$ against $H_1:\mu\neq\mu_0$.  The unrestricted MLE is
$\widehat\mu=\bar X$, and direct calculation gives
\[
    -2\log\Lambda
    =\frac{n(\bar X-\mu_0)^2}{\sigma^2}.
\]
Thus the level-$\alpha$ LRT rejects when
$\sqrt n\abs{\bar X-\mu_0}/\sigma>z_{1-\alpha/2}$.  Its finite-sample null
distribution is exactly standard normal after taking the signed square root;
no asymptotic approximation is needed.
\end{example}

\begin{example}[Two-sided coin test]
For $S\sim\Bin(n,p)$ and $H_0:p=p_0$ against $p\neq p_0$, the unrestricted
MLE is $\widehat p=S/n$.  The LRT statistic is
\[
 \Lambda(s)=
 \frac{p_0^s(1-p_0)^{n-s}}
      {(s/n)^s(1-s/n)^{n-s}},
\]
with $0^0=1$.  Rejection occurs in both tails, but discreteness means the two
tail probabilities generally cannot be balanced exactly without
randomization.
\end{example}

\section{One-Sided Hypotheses}

Now consider an ordered scalar parameter and hypotheses such as
\[
    H_0:\theta\leq\theta_0
    \qquad\text{versus}\qquad
    H_1:\theta>\theta_0.
\]
For many families, larger values of a statistic $T$ systematically favor
larger values of $\theta$.  This ordering can turn the collection of
simple-alternative Neyman--Pearson tests into a single optimal test.

\subsection{Uniformly Most Powerful (UMP) Tests}

\begin{definition}[UMP test]
A level-$\alpha$ test $\phi^*$ is \emph{uniformly most powerful} for
$H_0:\theta\in\Theta_0$ against $H_1:\theta\in\Theta_1$ if
\[
    \beta_{\phi^*}(\theta)\geq\beta_\phi(\theta)
    \quad\text{for every }\theta\in\Theta_1
\]
and every other level-$\alpha$ test $\phi$.
\end{definition}

The same test must be MP against every point in the alternative.  Such a test
often exists for one-sided problems because the likelihood-ratio ordering is
the same for all alternatives on one side.  It rarely exists for genuinely
two-sided problems: an MP test against a smaller parameter rejects in the
lower tail, while one against a larger parameter rejects in the upper tail.

For a boundary-null problem, the size is often attained at the boundary
$\theta_0$.  This must be proved, usually by showing that the rejection
probability is monotone in $\theta$.

\subsection{Monotone Likelihood Ratio (MLR)}

\begin{definition}[Monotone likelihood ratio]
A family $\{f_\theta:\theta\in\Theta\subseteq\R\}$ has a \emph{monotone
likelihood ratio} (MLR) in $T(X)$ if, whenever $\theta_1>\theta_0$, the ratio
$f_{\theta_1}(x)/f_{\theta_0}(x)$ is a nondecreasing function of $T(x)$ on the
common support.
\end{definition}

\begin{theorem}[Karlin--Rubin]
Suppose the family has MLR in $T$.  Choose $c$ and $\gamma$ so that
\[
    \phi^*(x)=\ind\{T(x)>c\}+\gamma\ind\{T(x)=c\}
\]
has rejection probability $\alpha$ at $\theta_0$.  Then $\phi^*$ is UMP
level $\alpha$ for
$H_0:\theta\leq\theta_0$ against $H_1:\theta>\theta_0$, provided its
rejection probability is nondecreasing in $\theta$.  Reversing the tails
gives the analogous result for a lower-sided alternative.
\end{theorem}

\begin{proof}[Proof sketch]
For every fixed $\theta_1>\theta_0$, the Neyman--Pearson lemma and MLR imply
that the same upper-tail test is MP for $\theta_0$ versus $\theta_1$.
Monotonicity of its power ensures level $\alpha$ throughout the composite
null.  Since it is MP at every point in the alternative, it is UMP.
\end{proof}

\begin{example}[One-sided weighted-coin test]
The binomial family has MLR in $S$ because the likelihood ratio displayed
above is increasing in $s$ whenever $p_1>p_0$.  Consequently, the test that
rejects for large $S$ is UMP level $\alpha$ for
$H_0:p\leq p_0$ against $H_1:p>p_0$.  Its exact power function is
\[
    \beta(p)=P_p(S>c)+\gamma P_p(S=c),
\]
which is nondecreasing in $p$ by stochastic ordering of binomial variables.
\end{example}

A regular one-parameter exponential family
\[
 f_\eta(x)=h(x)\exp\{\eta T(x)-A(\eta)\}
\]
has MLR in $T$, since for $\eta_1>\eta_0$ the likelihood ratio is a positive
constant times $e^{(\eta_1-\eta_0)T(x)}$.  Thus the theorem covers many
classical one-sided tests at once.

\section{Composite Hypotheses}

When both $\Theta_0$ and $\Theta_1$ contain multiple distributions, no single
ordering of the sample points need be optimal.  Useful approaches include
finding a UMP test when special structure permits it, restricting the class
of tests (for example to unbiased or invariant tests), using a likelihood
ratio test, or optimizing average or worst-case power.

The role of sufficiency persists.  If $T$ is sufficient and $\phi(X)$ is any
test, then
\[
    \phi^*(T)=\E_\theta[\phi(X)\mid T]
\]
is a well-defined, parameter-free randomized test with exactly the same power
function as $\phi$.  Thus attention may be restricted to tests based on a
sufficient statistic without loss.

\subsection{Uniformly Most Powerful (UMP) Tests}

UMP tests for composite alternatives exist only when the MP rejection regions
can be made compatible across all alternative points.  The one-sided MLR
setting is the principal example.  Two-sided alternatives usually require a
weaker optimality criterion.

\begin{proposition}[Nonexistence in a two-sided normal problem]
Let $X\sim\Normal(\mu,1)$ and test $H_0:\mu=0$ against $H_1:\mu\neq0$ at a
nontrivial level $\alpha$.  No UMP test exists.
\end{proposition}

\begin{proof}
By the Neyman--Pearson lemma, an MP test against any fixed $\mu_1>0$ rejects
for large $X$, whereas an MP test against any fixed $\mu_1<0$ rejects for
small $X$.  Apart from trivial levels, no test can have both forms.
\end{proof}

One response is to impose invariance under $x\mapsto-x$, leading to a
two-tailed test based on $\abs X$.  Another is to seek a UMP unbiased test.
Both restrictions rule out tests that spend nearly all type I error in one
tail.

\subsection{Least Favorable Priors}

A test can be derived as a Bayes rule by placing priors $\pi_0$ and $\pi_1$ on
the null and alternative.  The resulting mixture densities are
\[
    m_j(x)=\int_{\Theta_j}f_\theta(x)\pi_j(\dd\theta),
    \qquad j=0,1.
\]
For fixed prior odds and error costs, the Bayes test compares
$m_1(x)/m_0(x)$ with a constant.  Thus the Neyman--Pearson lemma applies to
the two mixture distributions even though the original hypotheses are
composite.

A least favorable prior places mass on parameter values for which control of
error or attainment of power is hardest.  If a Bayes test against the induced
mixtures has its null rejection probability bounded by $\alpha$ everywhere
and its worst alternative power equal to its average power under $\pi_1$,
then it attains a testing saddle point.  The Bayes lower bound and the
worst-case upper bound coincide, proving maximin optimality.

For a one-sided MLR problem, the point mass at the boundary $\theta_0$ is
often least favorable for size because the rejection probability is largest
there over the null.  With nuisance parameters or separated composite
hypotheses, least favorable priors can be nondegenerate mixtures and may only
exist as limiting sequences.

\section{Unbiased Tests}

The power of a level-$\alpha$ test can fall below $\alpha$ at some alternative
values.  Such a test rejects less often under those alternatives than under
the null boundary, an undesirable behavior that motivates a restriction.

\begin{definition}[Unbiased test]
A level-$\alpha$ test $\phi$ for $H_0:\theta\in\Theta_0$ against
$H_1:\theta\in\Theta_1$ is \emph{unbiased} if
\[
 \beta_\phi(\theta_1)\geq\beta_\phi(\theta_0)
 \quad\text{for every }\theta_1\in\Theta_1,
 \ \theta_0\in\Theta_0.
\]
In the common exact-size formulation this becomes
$\beta_\phi(\theta)\leq\alpha$ on $\Theta_0$ and
$\beta_\phi(\theta)\geq\alpha$ on $\Theta_1$.
\end{definition}

A test is \emph{uniformly most powerful unbiased} (UMPU) if it maximizes
power at every alternative point among all level-$\alpha$ unbiased tests.

Consider a regular one-parameter exponential family and the two-sided problem
$H_0:\eta=\eta_0$ against $H_1:\eta\neq\eta_0$.  If the power function is
differentiable, unbiasedness makes $\eta_0$ a local minimum, so
\begin{equation}
    \beta_\phi(\eta_0)=\alpha,
    \qquad \beta_\phi'(\eta_0)=0.
    \label{eq:unbiased-similarity-constraints}
\end{equation}
Differentiating under the integral and using the exponential-family score
gives
\[
 \beta_\phi'(\eta_0)
 =\E_{\eta_0}[\phi(X)\{T-\E_{\eta_0}T\}].
\]
Thus the constraints in \eqref{eq:unbiased-similarity-constraints} are
equivalent to
\[
    \E_{\eta_0}[\phi]=\alpha,
    \qquad
    \E_{\eta_0}[T\phi]=\alpha\E_{\eta_0}[T].
\]
The UMPU test rejects for $T<c_1$ or $T>c_2$, with possible boundary
randomization, where the two constants are chosen to satisfy these size and
moment constraints.  This allocates rejection probability between both tails
in the manner required by unbiasedness.

\begin{example}[Two-sided normal mean]
For $X_1,\ldots,X_n\iid\Normal(\mu,\sigma^2)$ with known $\sigma$, the UMPU
level-$\alpha$ test of $H_0:\mu=\mu_0$ against $\mu\neq\mu_0$ rejects when
\[
    \frac{\sqrt n\abs{\bar X-\mu_0}}{\sigma}>z_{1-\alpha/2}.
\]
Symmetry supplies the moment constraint, and the power is at least $\alpha$
for every $\mu\neq\mu_0$.
\end{example}

\subsection{Unbiased Tests with Nuisance Parameters}

Suppose the parameter is $(\psi,\lambda)$, where $\psi$ is of interest and
$\lambda$ is a nuisance parameter.  A level-$\alpha$ test must control
\[
    \sup_{\lambda}\E_{\psi_0,\lambda}[\phi(X)],
\]
which can be difficult if the null distribution depends on $\lambda$.

A test is \emph{similar} if
\begin{equation}
    \E_{\psi_0,\lambda}[\phi(X)]=\alpha
    \qquad\text{for every }\lambda.
    \label{eq:similar-test}
\end{equation}
Similarity is stronger than level control and is closely connected to
unbiasedness for two-sided problems.  Conditioning often removes the nuisance
parameter exactly.

Suppose that under $H_0$ a statistic $U$ is sufficient and complete for
$\lambda$.  If $\phi$ is similar, then
\[
 \E_{\psi_0,\lambda}
 \left[\E_{\psi_0,\lambda}(\phi\mid U)-\alpha\right]=0
 \quad\text{for every }\lambda.
\]
Completeness implies
\begin{equation}
    \E_{\psi_0,\lambda}[\phi(X)\mid U]=\alpha
    \quad\text{almost surely}.
    \label{eq:conditional-similarity}
\end{equation}
The conditional distribution used to construct the test is free of
$\lambda$, converting the problem into conditional simple-null tests.

\begin{example}[Normal mean with unknown variance]
Let $X_1,\ldots,X_n\iid\Normal(\mu,\sigma^2)$ and test
$H_0:\mu=\mu_0$ against $H_1:\mu\neq\mu_0$, with $\sigma^2$ unknown.  The
statistic
\[
    T=\frac{\sqrt n(\bar X-\mu_0)}{S_X}
\]
has a $t_{n-1}$ distribution under $H_0$, independent of $\sigma^2$.
Therefore
\[
    \phi(X)=\ind\{\abs T>t_{n-1,1-\alpha/2}\}
\]
is similar of level $\alpha$.  In the normal exponential family it is also
the UMP unbiased test for the two-sided mean problem.
\end{example}

\begin{example}[Conditional test for two binomial samples]
Let $S_j\sim\Bin(n_j,p_j)$ independently and test
$H_0:p_1=p_2=p$ against unequal success probabilities.  Under $H_0$, the
total $S_1+S_2$ is sufficient for the nuisance parameter $p$, and
\[
 P(S_1=s_1\mid S_1+S_2=s)
 =\frac{\binom{n_1}{s_1}\binom{n_2}{s-s_1}}
        {\binom{n_1+n_2}{s}},
\]
a hypergeometric probability independent of $p$.  Tail probabilities from
this conditional distribution yield Fisher's exact test.  Boundary
randomization can make it exactly similar; without randomization it is
typically conservative.
\end{example}

Conditioning is powerful but changes the criterion: optimality is established
within the similar or unbiased class, not necessarily among all level tests.
This qualification is important whenever a procedure is described as
``optimal'' in the presence of nuisance parameters.

\part{Optimal Estimation, Asymptotic}

\section{Introduction to Asymptotic Estimation}

Exact optimality compares estimators at a fixed sample size.  Such comparisons
can be difficult or overly restrictive: an unbiased estimator may not exist,
the finite-sample distribution may be intractable, and several procedures may
differ only by terms that vanish as the sample grows.  Asymptotic theory
studies a sequence of experiments and estimators
\[
    \delta_n=\delta_n(X_1,\ldots,X_n)
\]
as $n\to\infty$.

Three questions organize the theory:
\begin{enumerate}
    \item Does $\delta_n$ approach the target?  This is consistency.
    \item At what rate does its error vanish, and what is the limiting
    distribution of the rescaled error?
    \item Can another regular estimator have a smaller limiting risk or
    covariance?
\end{enumerate}

For regular finite-dimensional iid models, the canonical rate is
$n^{-1/2}$ and the canonical limit is Gaussian.  Neither conclusion is
automatic.  Boundary parameters, nonidentification, parameter-dependent
support, heavy tails, and nonsmooth criteria may produce different rates or
non-Gaussian limits.

\begin{example}[Weighted coin]
Let $X_1,\ldots,X_n\iid\Bern(p)$ and $\widehat p_n=\bar X_n$.  The law of
large numbers gives $\widehat p_n\convp p$, while the central limit theorem
gives, for $0<p<1$,
\[
    \sqrt n(\widehat p_n-p)
    \convd\Normal(0,p(1-p)).
\]
Consistency describes the eventual location of the estimator; asymptotic
normality describes the local shape and scale of its errors.
\end{example}

\section{Review}

All limits below are taken under $P_\theta$ unless another distribution is
specified.  Since the probability measure changes with $\theta$, asymptotic
claims should state whether they hold pointwise for each fixed $\theta$ or
uniformly over a subset of the parameter space.

\subsection{Convergence}

\begin{definition}[Modes of stochastic convergence]
For random variables $Y_n$ and $Y$:
\begin{itemize}
    \item $Y_n\as Y$ if $P(\lim_nY_n=Y)=1$;
    \item $Y_n\convp Y$ if, for every $\varepsilon>0$,
    $P(\abs{Y_n-Y}>\varepsilon)\to0$;
    \item $Y_n\convd Y$ if
    $\E[f(Y_n)]\to\E[f(Y)]$ for every bounded continuous $f$;
    \item $Y_n\to Y$ in $L^r$ if
    $\E\abs{Y_n-Y}^r\to0$.
\end{itemize}
The definitions extend to random vectors by replacing absolute value with a
norm.
\end{definition}

The usual implications are
\[
    L^r\Longrightarrow\text{probability}\Longrightarrow\text{distribution},
    \qquad
    \text{almost surely}\Longrightarrow\text{probability}.
\]
Convergence in distribution to a constant implies convergence in probability
to that constant.  In general, convergence in distribution does not imply
convergence of moments; uniform integrability is needed to control tails.

\begin{theorem}[Continuous mapping and Slutsky]
If $Y_n\convd Y$ and $g$ is continuous at $Y$ almost surely, then
$g(Y_n)\convd g(Y)$.  If additionally $Z_n\convp c$, then
\[
    Y_n+Z_n\convd Y+c,
    \qquad Y_nZ_n\convd cY,
\]
and $Y_n/Z_n\convd Y/c$ when $c\neq0$.
\end{theorem}

Stochastic order notation summarizes rates.  We write $Y_n=O_p(a_n)$ if
$Y_n/a_n$ is bounded in probability, and $Y_n=o_p(a_n)$ if
$Y_n/a_n\convp0$.  Thus
\[
    \sqrt n(\delta_n-\theta)=O_p(1)
    \quad\Longleftrightarrow\quad
    \delta_n-\theta=O_p(n^{-1/2}).
\]

\begin{theorem}[Weak law and central limit theorem]
If $X_1,X_2,\ldots$ are iid with mean $\mu$ and finite variance
$\sigma^2$, then
\[
    \bar X_n\convp\mu,
    \qquad
    \sqrt n(\bar X_n-\mu)\convd\Normal(0,\sigma^2).
\]
If only $\E\abs X_1<\infty$, the first conclusion still holds.
\end{theorem}

\subsection{Consistency}

\begin{definition}[Consistency]
An estimator $\delta_n$ of $g(\theta)$ is \emph{weakly consistent} at
$\theta$ if $\delta_n\convp g(\theta)$ under $P_\theta$.  It is strongly
consistent if $\delta_n\as g(\theta)$.  It is uniformly consistent on
$K\subseteq\Theta$ if, for every $\varepsilon>0$,
\[
    \sup_{\theta\in K}
    P_\theta\{\norm{\delta_n-g(\theta)}>\varepsilon\}\longrightarrow0.
\]
\end{definition}

Consistency alone gives neither a useful uncertainty scale nor good
finite-sample performance.  If $\delta_n$ is consistent, so is
$\delta_n+n^{-1/4}Z$ for many bounded random perturbations, even though the
perturbed estimator converges much more slowly than a root-$n$ estimator.

Two common proof strategies are direct moment bounds and optimization of a
sample criterion.  If $\E_\theta\delta_n\to g(\theta)$ and
$\Var_\theta(\delta_n)\to0$, Chebyshev's inequality proves consistency.  For
an $M$-estimator
\[
    \widehat\theta_n\in\argmax_{\vartheta\in\Theta}M_n(\vartheta),
\]
one instead shows that $M_n$ converges uniformly to a deterministic function
$M$ uniquely maximized at $\theta$.

\begin{theorem}[Argmax consistency]
Suppose $\Theta$ is compact,
\[
    \sup_{\vartheta\in\Theta}
    \abs{M_n(\vartheta)-M(\vartheta)}\convp0,
\]
and $M$ is uniquely and well-separatedly maximized at $\theta_0$: for every
neighborhood $U$ of $\theta_0$,
\[
    M(\theta_0)>\sup_{\vartheta\notin U}M(\vartheta).
\]
If $M_n(\widehat\theta_n)\geq\sup_\vartheta M_n(\vartheta)-o_p(1)$, then
$\widehat\theta_n\convp\theta_0$.
\end{theorem}

\begin{proof}[Proof sketch]
Outside any neighborhood $U$, the population criterion has a fixed gap below
its maximum.  Uniform convergence makes the sample criterion preserve this
gap with probability approaching one.  An approximate maximizer therefore
cannot lie outside $U$ with nonvanishing probability.
\end{proof}

\subsection{Asymptotic Normality}

\begin{definition}[Asymptotic normality]
An estimator is asymptotically normal with rate $r_n\to\infty$, asymptotic
bias $b(\theta)$, and asymptotic covariance $V(\theta)$ if
\[
    r_n\{\delta_n-g(\theta)\}
    \convd\Normal\{b(\theta),V(\theta)\}.
\]
In regular parametric models, $r_n=\sqrt n$.
\end{definition}

The most useful derivations start from an asymptotically linear expansion
\begin{equation}
    \delta_n-g(\theta)
    =\frac1n\sum_{i=1}^n\psi_\theta(X_i)+o_p(n^{-1/2}),
    \label{eq:asymptotic-linear-estimator}
\end{equation}
where $\E_\theta[\psi_\theta(X)]=0$ and
$\E_\theta\norm{\psi_\theta(X)}^2<\infty$.  The function $\psi_\theta$ is
the influence function.  The central limit theorem and Slutsky's theorem give
\[
    \sqrt n\{\delta_n-g(\theta)\}
    \convd\Normal\left(0,
       \E_\theta[\psi_\theta\psi_\theta^\mathsf T]\right).
\]

An estimated standard error $\widehat V_n^{1/2}/\sqrt n$ is useful only if
$\widehat V_n\convp V(\theta)$.  Studentization then gives
\[
    \frac{\sqrt n\{\delta_n-g(\theta)\}}{\sqrt{\widehat V_n}}
    \convd\Normal(0,1)
\]
in the scalar case, producing asymptotic confidence intervals.

\subsection{Delta Method}

\begin{theorem}[Delta method]
If
\[
    r_n(T_n-\theta)\convd Z
\]
and $g:\R^k\to\R^m$ is differentiable at $\theta$, then
\begin{equation}
    r_n\{g(T_n)-g(\theta)\}
    \convd \mathrm Dg(\theta)Z.
    \label{eq:delta-method}
\end{equation}
In particular, if $Z\sim\Normal_k(0,V)$, the limiting covariance is
$\mathrm Dg(\theta)V\mathrm Dg(\theta)^\mathsf T$.
\end{theorem}

\begin{proof}
Differentiability gives
\[
 g(T_n)-g(\theta)
 =\mathrm Dg(\theta)(T_n-\theta)+o(\norm{T_n-\theta}).
\]
Consistency and $r_n(T_n-\theta)=O_p(1)$ make the rescaled remainder $o_p(1)$.
Slutsky's theorem proves the result.
\end{proof}

\begin{example}[Log odds of a coin]
For $0<p<1$, let $g(p)=\log\{p/(1-p)\}$.  Since
$g'(p)=1/\{p(1-p)\}$ and
$\sqrt n(\widehat p-p)\convd\Normal(0,p(1-p))$,
\[
 \sqrt n\left[
 \log\frac{\widehat p}{1-\widehat p}
 -\log\frac p{1-p}\right]
 \convd\Normal\left(0,\frac1{p(1-p)}\right).
\]
This approximation is an interior result; the log odds is undefined when the
sample contains all heads or all tails and becomes unstable near $p=0$ or
$p=1$.
\end{example}

If the first derivative vanishes, the first-order delta method is degenerate.
If $g'(\theta)=0$ but $g''(\theta)\neq0$ and
$\sqrt n(T_n-\theta)\convd Z$, a second-order expansion yields
\[
    n\{g(T_n)-g(\theta)\}
    \convd\frac12g''(\theta)Z^2.
\]

\section{Asymptotic Unbiasedness}

There are several inequivalent definitions of asymptotic unbiasedness.  A
basic one requires
\begin{equation}
    \E_\theta[\delta_n]-g(\theta)\longrightarrow0.
    \label{eq:asymptotic-unbiasedness}
\end{equation}
Consistency does not by itself imply \eqref{eq:asymptotic-unbiasedness},
because rare but very large errors may prevent convergence of expectations.
Consistency plus uniform integrability of $\{\delta_n\}$ does imply it.

At the root-$n$ scale, one may instead require
\[
    \sqrt n\{\E_\theta[\delta_n]-g(\theta)\}\longrightarrow0.
\]
This stronger condition rules out first-order asymptotic bias.  If
\[
    \sqrt n\{\delta_n-g(\theta)\}
    \convd\Normal\{b(\theta),V(\theta)\}
\]
and the rescaled sequence is uniformly integrable, then its limiting scaled
bias is $b(\theta)$.

The distinction matters for mean squared error.  If moment convergence is
valid, then
\begin{equation}
 n\E_\theta[(\delta_n-g(\theta))^2]
 \longrightarrow V(\theta)+b(\theta)^2.
 \label{eq:asymptotic-mse}
\end{equation}
A small bias can reduce variance enough to improve total risk, just as in
finite-sample shrinkage.  Asymptotic unbiasedness is therefore a restriction,
not a universal requirement of good estimation.

\begin{example}[Smoothed coin estimator]
For constants $a,b>0$, the beta-prior posterior mean is
\[
    \widetilde p_n=\frac{S+a}{n+a+b}.
\]
Its bias is
\[
    \E_p[\widetilde p_n]-p
    =\frac{a-(a+b)p}{n+a+b}=O(n^{-1}),
\]
so it is unbiased to first order on the root-$n$ scale and has the same
asymptotic distribution as $S/n$.  Nevertheless, its finite-sample risk and
boundary behavior differ from those of the sample proportion.
\end{example}

\section{Asymptotic Efficiency}

If two consistent estimators have root-$n$ limits
\[
 \sqrt n(\delta_{j,n}-g(\theta))
 \convd\Normal(0,V_j(\theta)),\qquad j=1,2,
\]
then estimator $1$ is asymptotically at least as efficient as estimator $2$
when $V_2(\theta)-V_1(\theta)$ is positive semidefinite.  In the scalar case,
the asymptotic relative efficiency is often written
\[
    \operatorname{ARE}(1,2)=\frac{V_2(\theta)}{V_1(\theta)}.
\]
It approximates the factor by which estimator $2$ needs a larger sample to
match estimator $1$'s variance.

Pointwise asymptotic variance alone is too weak for a general optimality
theory.  An estimator can be modified at a single parameter value to converge
faster there while behaving normally elsewhere.

\begin{example}[Hodges' estimator]
Let $\bar X_n\sim\Normal(\theta,1/n)$ and define
\[
 \delta_n=
 \begin{cases}
  0,&\abs{\bar X_n}\leq n^{-1/4},\\
  \bar X_n,&\abs{\bar X_n}>n^{-1/4}.
 \end{cases}
\]
At $\theta=0$, $\sqrt n\delta_n\convp0$, apparently improving on the usual
$\Normal(0,1)$ limit.  At every fixed $\theta\neq0$, it has the same
asymptotic distribution as $\bar X_n$.  But along parameters approaching zero
near the threshold, its risk is poor.  This \emph{superefficiency} shows why
efficiency bounds require regularity under local alternatives, not merely a
pointwise limit at each fixed parameter.
\end{example}

\begin{definition}[Regular estimator]
In a locally asymptotically normal model, an estimator $\delta_n$ of a scalar
or vector parameter is regular at $\theta$ if, for every fixed local direction
$h$, the distribution of
\[
    \sqrt n\{\delta_n-(\theta+h/\sqrt n)\}
\]
under $P_{\theta+h/\sqrt n}$ converges to a limit that does not depend on $h$.
\end{definition}

Regularity excludes the unstable local behavior of Hodges' estimator and is
the appropriate class for asymptotic information bounds.

\section{Asymptotic Information Inequality}

Let $s_\theta(X)=\grad_\theta\log f_\theta(X)$ and let
$I(\theta)=\E_\theta[s_\theta s_\theta^\mathsf T]$.  For a regular estimator
of $g(\theta)\in\R^m$, the asymptotic analogue of the Cram\'er--Rao bound is
\begin{equation}
    V(\theta)\succeq
    \mathrm Dg(\theta)I(\theta)^{-1}
    \mathrm Dg(\theta)^\mathsf T.
    \label{eq:asymptotic-information-bound}
\end{equation}
Unlike the finite-sample bound, this result allows small finite-sample bias;
regularity supplies the local stability that replaces exact unbiasedness.

\begin{theorem}[Hájek--Le Cam convolution theorem]
In a regular LAN model, the limiting distribution of a regular estimator of
$\theta$ has the representation
\[
    Z=Z_0+W,
    \qquad
    Z_0\sim\Normal_k(0,I(\theta)^{-1}),
\]
where $Z_0$ and $W$ are independent.  Consequently
$\Cov(Z)\succeq I(\theta)^{-1}$ whenever the covariance exists.
\end{theorem}

The Gaussian component $Z_0$ is unavoidable sampling noise; $W$ is additional
noise contributed by an inefficient procedure.  An estimator is
\emph{asymptotically efficient} if $W=0$, so that
\[
    \sqrt n(\delta_n-\theta)
    \convd\Normal_k(0,I(\theta)^{-1}).
\]

The corresponding local asymptotic minimax theorem provides a risk statement
that is robust to moving parameters.  For suitable bowl-shaped loss $L$ and
every fixed radius $c$,
\[
 \liminf_{n\to\infty}
 \sup_{\norm h\leq c}
 \E_{\theta+h/\sqrt n}
 L\!\left(\sqrt n\{\delta_n-(\theta+h/\sqrt n)\}\right)
\]
is bounded below, as $c\to\infty$, by the risk of
$Z_0\sim\Normal_k(0,I(\theta)^{-1})$.  This formulation rules out gains at a
single point purchased by poor behavior in a shrinking neighborhood.

\begin{example}[Efficiency for the coin bias]
One Bernoulli observation has information $I(p)=1/\{p(1-p)\}$.  Hence the
asymptotic variance lower bound for regular estimators of $p$ is
$I(p)^{-1}=p(1-p)$.  Since
\[
    \sqrt n(\widehat p-p)\convd\Normal(0,p(1-p)),
\]
the sample proportion is asymptotically efficient at every interior
$p\in(0,1)$.  The regular theory does not apply at $p=0$ or $p=1$.
\end{example}

\section{Maximum Likelihood Estimators}

For a dominated model, an MLE is any measurable maximizer
\[
    \widehat\theta_n\in\argmax_{\theta\in\Theta}
       \ell_n(\theta),
    \qquad
    \ell_n(\theta)=\sum_{i=1}^n\log f_\theta(X_i).
\]
The MLE is invariant under one-to-one reparameterization: if
$\widehat\theta_n$ maximizes the likelihood and $g$ is one-to-one, then
$g(\widehat\theta_n)$ is the MLE of $g(\theta)$.  For non-one-to-one $g$, the
same conclusion holds when the likelihood for $g(\theta)$ is defined by
profiling over its fibers.

\subsection{Consistency and Inconsistency}

If the data come from $P_{\theta_0}$, the normalized log-likelihood tends
pointwise to
\[
    M(\theta)=\E_{\theta_0}[\log f_\theta(X)].
\]
The difference from its value at the truth is
\begin{equation}
    M(\theta)-M(\theta_0)
    =-\operatorname{KL}(P_{\theta_0}\Vert P_\theta)\leq0.
    \label{eq:mle-kl-identity}
\end{equation}
Thus identifiability makes $\theta_0$ the unique population maximizer.  To
transfer this fact to the sample maximizer, pointwise laws of large numbers
are generally insufficient; one needs uniform control and must prevent the
MLE from escaping to the boundary or infinity.

\begin{theorem}[A consistency theorem for MLEs]
Suppose $\Theta$ is compact, the model is identifiable, the map
$\theta\mapsto\log f_\theta(x)$ is suitably continuous, and a uniform law of
large numbers holds:
\[
 \sup_{\theta\in\Theta}
 \abs{n^{-1}\ell_n(\theta)-M(\theta)}\convp0.
\]
If $M$ is well-separatedly maximized at $\theta_0$, then every approximate MLE
is consistent for $\theta_0$.
\end{theorem}

\begin{example}[Coin MLE]
For $S\sim\Bin(n,p)$, the log-likelihood is
\[
    \ell_n(p)=S\log p+(n-S)\log(1-p),
\]
and the MLE is $\widehat p=S/n$, including the boundary values when all flips
agree.  The law of large numbers immediately proves consistency.
\end{example}

MLEs may be inconsistent when the model is misspecified, the parameter is not
identified, the dimension grows too quickly, or maximization exploits an
unbounded likelihood.

\begin{example}[Misspecification]
If the true distribution is $P_0$ but the fitted family
$\{P_\theta\}$ does not contain it, the MLE generally converges to the
pseudo-true parameter
\[
    \theta^*=\argmin_\theta
       \operatorname{KL}(P_0\Vert P_\theta),
\]
not to a literal data-generating parameter.  Inference then uses a sandwich
covariance rather than the inverse Fisher information unless the model is
correctly specified.
\end{example}

\begin{example}[An unbounded likelihood]
In a normal mixture with an unknown component variance, a component mean can
be placed on one observation while its variance tends to zero.  The
likelihood then diverges, so an unconstrained global MLE does not exist.
Large sample size does not repair a nonexistent maximizer; constraints or
penalization are required.
\end{example}

\subsection{Asymptotic Normality and Efficiency}

The score equation is $\grad\ell_n(\widehat\theta_n)=0$.  For a scalar
interior parameter, Taylor expansion around the truth gives
\[
 0=\ell_n'(\theta_0)
   +\ell_n''(\widetilde\theta_n)(\widehat\theta_n-\theta_0),
\]
where $\widetilde\theta_n$ lies between $\widehat\theta_n$ and $\theta_0$.
Therefore
\begin{equation}
 \sqrt n(\widehat\theta_n-\theta_0)
 =\left\{-\frac1n\ell_n''(\widetilde\theta_n)\right\}^{-1}
   \frac{\ell_n'(\theta_0)}{\sqrt n}.
 \label{eq:mle-score-expansion}
\end{equation}
Under regularity and consistency, the first factor converges to
$I(\theta_0)^{-1}$ and the second converges to
$\Normal(0,I(\theta_0))$.  Hence
\begin{equation}
    \sqrt n(\widehat\theta_n-\theta_0)
    \convd\Normal(0,I(\theta_0)^{-1}).
    \label{eq:mle-asymptotic-normality}
\end{equation}
The MLE attains the asymptotic information bound and is efficient.

Equation \eqref{eq:mle-score-expansion} also gives the asymptotically linear
representation
\[
 \widehat\theta_n-\theta_0
 =\frac1n\sum_{i=1}^n
 I(\theta_0)^{-1}s_{\theta_0}(X_i)+o_p(n^{-1/2}).
\]
Thus the efficient influence function is
$I(\theta_0)^{-1}s_{\theta_0}(X)$.

Common variance estimators are the inverse expected information, inverse
observed information, and the outer product of scores.  Under correct
specification they are asymptotically equivalent.  Under misspecification,
if
\[
 A=-\E_0[\Hess_\theta\log f_{\theta^*}(X)],
 \qquad
 B=\E_0[s_{\theta^*}(X)s_{\theta^*}(X)^\mathsf T],
\]
then the limiting covariance is the sandwich matrix
\begin{equation}
    A^{-1}BA^{-1},
    \label{eq:sandwich-covariance}
\end{equation}
which reduces to $I(\theta_0)^{-1}$ when the information identity $A=B$
holds.

\subsection{MLE for Multi-parameter Families}

Let $\theta\in\R^k$.  Under the multivariate analogues of consistency,
interiority, differentiability, nonsingular information, and uniform control
of the Hessian,
\begin{equation}
 \sqrt n(\widehat\theta_n-\theta_0)
 =I(\theta_0)^{-1}
   \frac1{\sqrt n}\sum_{i=1}^n s_{\theta_0}(X_i)+o_p(1)
 \convd\Normal_k(0,I(\theta_0)^{-1}).
 \label{eq:multivariate-mle-limit}
\end{equation}
For a differentiable target $g:\R^k\to\R^m$, the delta method gives
\[
 \sqrt n\{g(\widehat\theta_n)-g(\theta_0)\}
 \convd\Normal_m\left(0,
 \mathrm Dg(\theta_0)I(\theta_0)^{-1}
 \mathrm Dg(\theta_0)^\mathsf T\right).
\]

\begin{example}[Mean and variance of a normal distribution]
If $X_i\iid\Normal(\mu,\sigma^2)$, parameterize by
$\theta=(\mu,\sigma^2)$.  The MLEs are
\[
    \widehat\mu=\bar X,
    \qquad
    \widehat\sigma^2=\frac1n\sum_{i=1}^n(X_i-\bar X)^2.
\]
The per-observation information matrix is diagonal:
\[
 I(\mu,\sigma^2)=
 \begin{pmatrix}
  1/\sigma^2&0\\
  0&1/(2\sigma^4)
 \end{pmatrix}.
\]
Consequently,
\[
 \sqrt n
 \begin{pmatrix}
  \widehat\mu-\mu\\
  \widehat\sigma^2-\sigma^2
 \end{pmatrix}
 \convd\Normal_2\left(
 \begin{pmatrix}0\\0\end{pmatrix},
 \begin{pmatrix}
  \sigma^2&0\\
  0&2\sigma^4
 \end{pmatrix}\right).
\]
Although $\widehat\sigma^2$ is biased at finite $n$, its bias is $O(n^{-1})$
and disappears on the root-$n$ scale.
\end{example}

If $\theta=(\psi,\lambda)$ contains a nuisance parameter, the asymptotic
covariance for estimating $\psi$ is the corresponding block of
$I(\theta)^{-1}$, equivalently the inverse Schur complement
\[
 \left(I_{\psi\psi}
 -I_{\psi\lambda}I_{\lambda\lambda}^{-1}
  I_{\lambda\psi}\right)^{-1}.
\]
This is the same efficient information that governs local asymptotic testing.
Estimation and testing are therefore two views of the same local Gaussian
geometry.

Finally, \eqref{eq:multivariate-mle-limit} is an interior regular-model
result.  At boundaries, under singular information, or when the number of
parameters grows with $n$, the MLE may have a different rate and limit.
Those cases require a fresh expansion rather than mechanical use of inverse
Fisher information.

\part{Optimal Hypothesis Testing, Asymptotic}

Exact finite-sample optimality is available only in specially structured
models.  In regular problems, asymptotic theory replaces the experiment based
on $n$ observations by a limiting Gaussian experiment.  This approximation
does more than provide critical values: it identifies the scale of
statistically distinguishable alternatives and yields upper bounds on local
power.

\section{Asymptotic Testing Framework}

Let $X^{(n)}$ denote the data in the $n$th experiment, with distribution
$P_{n,\theta}$, and let $\phi_n(X^{(n)})\in[0,1]$ be a sequence of tests.  The
models need not literally be iid, although that is the principal example.

\begin{definition}[Asymptotic level]
The sequence $\{\phi_n\}$ has \emph{pointwise asymptotic level $\alpha$} if
\[
    \limsup_{n\to\infty}\E_{n,\theta}[\phi_n]\leq\alpha
    \qquad\text{for every }\theta\in\Theta_0.
\]
It has \emph{uniform asymptotic level $\alpha$} if
\begin{equation}
    \limsup_{n\to\infty}
    \sup_{\theta\in\Theta_0}\E_{n,\theta}[\phi_n]
    \leq\alpha.
    \label{eq:uniform-asymptotic-level}
\end{equation}
\end{definition}

Uniform validity is stronger: the parameter value producing the largest
rejection probability may vary with $n$.  Pointwise convergence alone does
not exclude severe size distortion along such a sequence.

\begin{definition}[Consistency of a test]
A sequence of tests is \emph{consistent} against a fixed
$\theta\in\Theta_1$ if
\[
    \E_{n,\theta}[\phi_n]\longrightarrow1.
\]
It is uniformly consistent over a set $A\subseteq\Theta_1$ if the infimum of
its power over $A$ converges to one.
\end{definition}

Consistency against fixed alternatives is often easy and does not distinguish
good tests from mediocre ones.  More informative comparisons use local
alternatives that approach the null as the sample grows.

\subsection{Fixed and Local Alternatives}

In a regular iid model, estimators fluctuate on the $n^{-1/2}$ scale.  Thus
the canonical local alternatives to a point null $\theta_0$ are
\begin{equation}
    \theta_n=\theta_0+\frac{h}{\sqrt n},
    \qquad h\in\R^k\text{ fixed}.
    \label{eq:local-alternatives}
\end{equation}
Alternatives separated from $\theta_0$ by much more than $n^{-1/2}$ are
typically detected with probability tending to one.  Those separated by much
less are typically indistinguishable from the null.  At the $n^{-1/2}$ scale,
power has a nondegenerate limit strictly between size and one.

\begin{example}[Local alternatives for a weighted coin]
Let $X_1,\ldots,X_n\iid\Bern(p)$ and test $H_0:p=p_0$, where
$0<p_0<1$.  Define
\[
    Z_n=\frac{S-np_0}{\sqrt{np_0(1-p_0)}}.
\]
Under $p=p_0$, the central limit theorem gives $Z_n\convd\Normal(0,1)$.
Under $p_n=p_0+h/\sqrt n$,
\[
 Z_n
 =\frac{S-np_n}{\sqrt{np_n(1-p_n)}}
   \sqrt{\frac{p_n(1-p_n)}{p_0(1-p_0)}}
   +\frac{h}{\sqrt{p_0(1-p_0)}},
\]
so
\begin{equation}
    Z_n\convd
    \Normal\left(\frac{h}{\sqrt{p_0(1-p_0)}},1\right).
    \label{eq:coin-local-normal}
\end{equation}
The upper-tailed test $\ind\{Z_n>z_{1-\alpha}\}$ therefore has limiting
local power
\[
    1-\Phi\left(z_{1-\alpha}
       -\frac{h}{\sqrt{p_0(1-p_0)}}\right).
\]
\end{example}

\subsection{Asymptotic Critical Values and p-Values}

Suppose a statistic $T_n$ has continuous limiting distribution $F_0$ under
the null.  If large values are evidence against $H_0$, then
\[
    \phi_n=\ind\{T_n>F_0^{-1}(1-\alpha)\}
\]
has asymptotic rejection probability $\alpha$ at parameter values for which
$T_n\convd F_0$.  The corresponding asymptotic $p$-value is
$1-F_0(T_n)$.

This argument proves pointwise validity only.  Uniform validity requires a
uniform distributional approximation over the composite null.  It can fail
near boundaries, where information degenerates, or when nuisance parameters
are weakly identified.  When a finite-sample null distribution is available,
there is no reason to replace it with an asymptotic one.

\section{Local Asymptotic Normality}

Let $X_1,\ldots,X_n\iid P_\theta$ with log-likelihood
$\ell_n(\theta)=\sum_{i=1}^n\log f_\theta(X_i)$.  Write
\[
    s_\theta(X)=\grad_\theta\log f_\theta(X),
    \qquad
    I(\theta)=\E_\theta[s_\theta(X)s_\theta(X)^\mathsf T].
\]
The normalized score at $\theta_0$ is
\[
    \Delta_n=\frac{1}{\sqrt n}\sum_{i=1}^n s_{\theta_0}(X_i).
\]
Under regularity conditions, the multivariate central limit theorem gives
$\Delta_n\convd\Normal_k(0,I(\theta_0))$ under $P_{\theta_0}$.

\begin{definition}[Local asymptotic normality]
The sequence of experiments is \emph{locally asymptotically normal} (LAN) at
$\theta_0$ if, for every fixed $h\in\R^k$,
\begin{equation}
 \log\frac{\dd P_{n,\theta_0+h/\sqrt n}}
              {\dd P_{n,\theta_0}}
 =h^\mathsf T\Delta_n
  -\frac12h^\mathsf TI(\theta_0)h+o_{P_{\theta_0}}(1),
 \label{eq:lan-expansion}
\end{equation}
with $\Delta_n\convd\Normal_k(0,I(\theta_0))$ under the null.
\end{definition}

A second-order Taylor expansion of the log-likelihood suggests
\eqref{eq:lan-expansion}; the LAN statement additionally controls the
remainder under local perturbations.  The limiting log-likelihood ratio is
exactly that of observing
\[
    Y\sim\Normal_k(I(\theta_0)h,I(\theta_0))
\]
with $h=0$ under the null.  Hence regular local testing problems inherit the
geometry of a Gaussian shift experiment.

\subsection{Contiguity}

\begin{definition}[Contiguity]
A sequence $Q_n$ is \emph{contiguous} with respect to $P_n$, written
$Q_n\triangleleft P_n$, if $P_n(A_n)\to0$ implies $Q_n(A_n)\to0$ for every
sequence of events $A_n$.  If contiguity holds in both directions, the
sequences are mutually contiguous.
\end{definition}

In a regular LAN model, $P_{n,\theta_0+h/\sqrt n}$ and
$P_{n,\theta_0}$ are mutually contiguous.  Thus an event that is asymptotically
negligible under the null remains negligible under local alternatives.  This
is what allows stochastic expansions proved under the null to remain useful
for local-power calculations.

\begin{theorem}[Le Cam's third lemma, Gaussian form]
Suppose under $P_n$,
\[
 \left(T_n,\log\frac{\dd Q_n}{\dd P_n}\right)
 \convd \Normal\left(
 \begin{pmatrix}\mu\\-\sigma^2/2\end{pmatrix},
 \begin{pmatrix}V&c\\c^\mathsf T&\sigma^2\end{pmatrix}
 \right).
\]
Then under $Q_n$, $T_n\convd\Normal(\mu+c,V)$, with the evident vector
interpretation.
\end{theorem}

The lemma says that exponential tilting by the local likelihood ratio shifts
the limiting mean by the covariance with that likelihood ratio.  Applied to
LAN, it yields
\begin{equation}
    \Delta_n\convd
    \Normal_k(I(\theta_0)h,I(\theta_0))
    \quad\text{under }\theta_0+h/\sqrt n.
    \label{eq:score-under-local}
\end{equation}

\subsection{Asymptotic Power Envelope}

For a scalar parameter, standardize the central sequence as
$Z_n=I(\theta_0)^{-1/2}\Delta_n$.  Under the null it converges to
$\Normal(0,1)$; under $\theta_0+h/\sqrt n$ it converges to
$\Normal(h\sqrt{I(\theta_0)},1)$.  The Neyman--Pearson lemma in the limiting
Gaussian experiment gives the one-sided power envelope
\begin{equation}
    \pi^*(h)=1-\Phi\{z_{1-\alpha}-h\sqrt{I(\theta_0)}\},
    \qquad h>0.
    \label{eq:local-power-envelope}
\end{equation}
A test whose limiting local power equals this expression is locally
asymptotically most powerful in the specified direction.

For two-sided testing, no UMP test exists even in the limiting Gaussian
experiment.  Restricting to asymptotically unbiased or symmetric tests gives
the rejection rule $\abs{Z_n}>z_{1-\alpha/2}$ and local power
\[
 P\left(\abs{Z+h\sqrt{I(\theta_0)}}>z_{1-\alpha/2}\right),
 \qquad Z\sim\Normal(0,1).
\]

\section{Likelihood-Based Tests}

Let $\widehat\theta$ be the unrestricted maximum likelihood estimator and
$\widetilde\theta$ the MLE under $H_0$.  Three classical tests measure the
same local departure from the null in different ways:
\begin{itemize}
    \item the likelihood-ratio test compares maximized likelihoods;
    \item the Wald test measures the distance from $\widehat\theta$ to the
    null using estimated standard errors;
    \item the score test measures the slope of the likelihood at the
    constrained estimate $\widetilde\theta$.
\end{itemize}

\subsection{Likelihood-Ratio Tests and Wilks' Theorem}

Define
\begin{equation}
    W_n=2\{\ell_n(\widehat\theta)-\ell_n(\widetilde\theta)\}
       =-2\log\Lambda_n.
    \label{eq:asymptotic-lr-statistic}
\end{equation}

\begin{theorem}[Wilks]
Suppose the model is regular, the true parameter is an interior point of
$\Theta_0$, and $\Theta_0$ is a smooth submanifold of codimension $r$.  Under
$H_0$,
\[
    W_n\convd\chi_r^2.
\]
Thus the asymptotic level-$\alpha$ LRT rejects when
$W_n>\chi^2_{r,1-\alpha}$.
\end{theorem}

The degrees of freedom equal the number of independent restrictions imposed
by the null.  A quadratic expansion of $\ell_n$ shows that $W_n$ is
asymptotically the squared length of the normalized score component
orthogonal to the null tangent space.

Under local alternatives, the limit is noncentral chi-squared:
\begin{equation}
    W_n\convd\chi_r^2(\lambda),
    \label{eq:noncentral-wilks}
\end{equation}
where $\lambda$ is the squared information distance from the local direction
$h$ to the tangent space of the null.  This limit determines local power.

\subsection{Wald Tests}

For a scalar restriction $H_0:g(\theta)=0$, the Wald statistic is
\begin{equation}
    W_n^{\mathrm{Wald}}
    =\frac{g(\widehat\theta)^2}
    {\grad g(\widehat\theta)^\mathsf T
     \{nI(\widehat\theta)\}^{-1}\grad g(\widehat\theta)}.
    \label{eq:wald-statistic}
\end{equation}
Under regularity and the null, it converges to $\chi_1^2$.  For $r$ smooth
restrictions, replace the denominator by the corresponding covariance matrix
and obtain a $\chi_r^2$ limit.

Wald tests are easy to construct from an estimator and its standard error,
but their finite-sample behavior is not invariant to nonlinear
reparameterization of the restriction.  They can also behave poorly when the
MLE is near a boundary or its estimated variance is unstable.

\subsection{Score Tests}

The score, or Lagrange multiplier, test needs only the null fit.  For a scalar
parameter and a simple null, its statistic is
\begin{equation}
    W_n^{\mathrm{score}}
    =\frac{\{\ell_n'(\theta_0)\}^2}{nI(\theta_0)},
    \label{eq:score-test-statistic}
\end{equation}
which converges to $\chi_1^2$ under $H_0$.  A signed one-sided version uses
$\ell_n'(\theta_0)/\sqrt{nI(\theta_0)}$ directly.

\begin{example}[Three asymptotic coin tests]
For $S\sim\Bin(n,p)$ and $H_0:p=p_0$, the signed score statistic is
\[
    Z_{\mathrm{score}}
    =\frac{S-np_0}{\sqrt{np_0(1-p_0)}}.
\]
The Wald statistic instead standardizes $\widehat p-p_0$ using
$\widehat p(1-\widehat p)/n$, where $\widehat p=S/n$.  The likelihood-ratio
statistic is
\[
 2\left[S\log\frac{\widehat p}{p_0}
 +(n-S)\log\frac{1-\widehat p}{1-p_0}\right].
\]
For fixed interior $p_0$, all three have the same first-order null limit and
the same local asymptotic power.  They can differ noticeably for small $n$ or
when $p_0$ is close to zero or one.
\end{example}

\subsection{First-Order Equivalence}

In a regular model with a scalar parameter,
\[
 2\{\ell_n(\widehat\theta)-\ell_n(\theta_0)\}
 =nI(\theta_0)(\widehat\theta-\theta_0)^2+o_p(1)
 =\frac{\ell_n'(\theta_0)^2}{nI(\theta_0)}+o_p(1).
\]
Consequently the likelihood-ratio, Wald, and score tests are first-order
equivalent under the null and contiguous alternatives.  This equivalence does
not imply equal finite-sample performance or equal higher-order accuracy.

\section{Nuisance Parameters and Efficient Scores}

Partition the parameter as $\theta=(\psi,\lambda)$, where
$\psi\in\R^r$ is the parameter of interest and $\lambda$ is nuisance.  Write
the score and information in blocks:
\[
 s_\theta=
 \begin{pmatrix}s_\psi\\s_\lambda\end{pmatrix},
 \qquad
 I(\theta)=
 \begin{pmatrix}
 I_{\psi\psi}&I_{\psi\lambda}\\
 I_{\lambda\psi}&I_{\lambda\lambda}
 \end{pmatrix}.
\]

\subsection{Efficient Score and Information}

The component of $s_\psi$ linearly predictable from the nuisance score does
not provide information that distinguishes a change in $\psi$ from a change
in $\lambda$.  Removing it gives the efficient score
\begin{equation}
    s_{\psi\cdot\lambda}
    =s_\psi-I_{\psi\lambda}I_{\lambda\lambda}^{-1}s_\lambda.
    \label{eq:efficient-score}
\end{equation}
Its covariance is the efficient information
\begin{equation}
    I_{\psi\cdot\lambda}
    =I_{\psi\psi}
     -I_{\psi\lambda}I_{\lambda\lambda}^{-1}I_{\lambda\psi},
    \label{eq:efficient-information}
\end{equation}
the Schur complement of $I_{\lambda\lambda}$.  Score tests with nuisance
parameters use the efficient score evaluated at the restricted estimator.

Geometrically, nuisance scores span a tangent space in $L^2(P_\theta)$.
Equation \eqref{eq:efficient-score} projects the score for the parameter of
interest onto the orthogonal complement of that space.  Local power bounds
are obtained by replacing ordinary information with efficient information.

\subsection{Plug-In and Profile Likelihood}

The profile log-likelihood is
\[
    \ell_n^p(\psi)=\sup_\lambda\ell_n(\psi,\lambda).
\]
Comparing its constrained and unconstrained maxima yields the usual
likelihood-ratio statistic.  Under regularity, estimating the nuisance
parameter changes the first-order variance precisely through
$I_{\psi\cdot\lambda}$, and Wilks' theorem still gives degrees of freedom equal
to the dimension of the restriction on $\psi$.

Plug-in critical values are valid when the estimated nuisance parameter is
consistent and the null approximation is sufficiently uniform.  Exact
pivots, such as the Student $t$ statistic, are preferable when available.
Otherwise one may use a parametric bootstrap, but its validity is itself an
asymptotic theorem and can fail in nonregular problems.

\section{Asymptotic Optimality}

An asymptotically valid test is not automatically optimal.  In LAN models,
optimality is defined by comparison with the Gaussian limit experiment.

\subsection{Locally Asymptotically Most Powerful Tests}

For testing a scalar $H_0:\theta=\theta_0$ against local upper alternatives,
the normalized score test
\[
    \phi_n=\ind\left\{
      I(\theta_0)^{-1/2}\Delta_n>z_{1-\alpha}
    \right\}
\]
attains the envelope \eqref{eq:local-power-envelope}.  No asymptotic
level-$\alpha$ test satisfying the regularity conditions can have larger
limiting power at the same local direction $h>0$.  Likelihood-ratio and Wald
tests have the same first-order local power when their rejection regions are
oriented in the same one-sided direction.

With a vector parameter, no test is most powerful in every direction.  For a
specified direction $h$, the optimal statistic is proportional to
$h^\mathsf T\Delta_n$.  For an unspecified departure from a smooth null,
the likelihood-ratio statistic uses the squared norm of the efficient score
in the normal space to the null.  This yields rotationally invariant, rather
than directionwise, optimality in the limiting Gaussian experiment.

\subsection{Pitman Efficiency}

Suppose two asymptotic level-$\alpha$ tests have limiting local powers of the
form
\[
    1-\Phi(z_{1-\alpha}-c_jh),\qquad j=1,2.
\]
Their Pitman asymptotic relative efficiency is
\[
    \operatorname{ARE}(1,2)=\frac{c_1^2}{c_2^2}.
\]
It is interpreted as the limiting ratio of sample sizes required for equal
power against nearby alternatives.  ARE compares first-order local behavior;
it does not determine performance against fixed alternatives or in small
samples.

\subsection{Testing and Confidence Sets}

A family of level-$\alpha$ tests can be inverted to form a confidence set:
\[
    C_n(X)=\{\theta_0:\text{the test of }H_0:\theta=\theta_0
                   \text{ does not reject}\}.
\]
If the tests have uniform asymptotic level $\alpha$, then
\[
    \liminf_{n\to\infty}\inf_{\theta\in\Theta}
    P_\theta\{\theta\in C_n(X)\}\geq1-\alpha.
\]
Pointwise level yields only pointwise coverage.  Wald, score, and
likelihood-ratio confidence regions arise by inverting their corresponding
tests; despite first-order equivalence, their finite-sample shapes can be
quite different.

\section{Nonregular Problems}

The standard theory relies on an interior true parameter, differentiability,
nonsingular information, identifiable parameters, and a stable support.
When these conditions fail, the $n^{-1/2}$ rate, LAN expansion, or chi-squared
calibration may be wrong.

\subsection{Parameters on the Boundary}

Suppose a scalar parameter is constrained to $\theta\geq0$ and the null is
$H_0:\theta=0$.  The null lies on the boundary, so the unrestricted local
Gaussian approximation must itself be projected onto a half-line.  In a
standard one-dimensional case,
\[
    -2\log\Lambda_n
    \convd \tfrac12\chi_0^2+\tfrac12\chi_1^2,
\]
where $\chi_0^2$ is a point mass at zero.  Using an ordinary $\chi_1^2$
critical value is then conservative.  More complicated inequality
constraints lead to chi-bar-squared mixtures whose weights depend on the
local geometry.

The coin model illustrates another boundary failure.  At $p_0=0$, every
observation is zero under the null and the Fisher information formula
$1/\{p_0(1-p_0)\}$ is not meaningful.  Alternatives on the scale $1/n$, not
$1/\sqrt n$, can be detectable because
\[
 P_{p_n}(S=0)=(1-p_n)^n\longrightarrow e^{-c}
 \quad\text{when }p_n=c/n.
\]
The limiting experiment is Poisson rather than Gaussian.

\subsection{Nonidentification and Singular Information}

If a nuisance parameter is present only under the alternative, it cannot be
consistently estimated under the null.  Mixture-component tests and tests for
an unknown change point are common examples.  The likelihood-ratio statistic
may converge to the supremum of a stochastic process rather than a
chi-squared variable.

Similarly, singular Fisher information signals that the quadratic
log-likelihood approximation is inadequate.  Higher-order terms may determine
the rate and limiting distribution.  These cases require problem-specific
analysis; reporting a standard Wald or Wilks $p$-value is not justified merely
because the sample size is large.

\subsection{Resampling Calibration}

A parametric bootstrap fits the null model, simulates data from the fitted
null, and recomputes the statistic to approximate its null distribution.  A
permutation test instead exploits exchangeability under the null and may be
finite-sample exact.  The distinction is important:
\begin{itemize}
    \item permutation validity follows from a group symmetry and can be exact;
    \item bootstrap validity follows from consistency of an estimated
    distribution and is generally asymptotic;
    \item neither method automatically repairs nonidentification or a
    discontinuous statistic.
\end{itemize}

The practical asymptotic workflow is therefore: identify the null geometry,
derive the statistic's limit under the null and contiguous alternatives,
verify whether convergence is uniform over nuisance parameters, and only then
choose analytical or resampling critical values.  Local power can be compared
with the LAN Gaussian envelope when regularity holds; when it does not, the
appropriate non-Gaussian limiting experiment determines what optimality can
mean.

\end{document}
